%
%
%
%
%
%
%
%
%
%
%
%
%
%
%
%
%
%
%
%
%
%

% ── Platz: Dichte, Zeilenhöhe, Umbruchreserve ────────────────────────
% Müssen VOR styles/30_layout_spacing stehen.
% Dichte und Zeilenhöhe sind getrennt, weil unterschiedlich riskant:
% Abstände sind Leerraum und vertragen jede Skalierung, Zeilenhöhe
% enthält die Schrift selbst — nur in kleinen Schritten senken und
% danach das PDF auf kollidierende Formelzeilen prüfen.
% \def\ZSFDensityFactor{0.85}     % 1    = Referenz, kleiner = dichter
% \def\ZSFLeadingFactor{0.92}     % 0.95 = Referenz, kleiner = enger
%
% Bereichsfaktoren multiplizieren den globalen Faktor für IHREN Bereich —
% für ZSF, die ungleich verteilt sind (box-lastig vs. textlastig).
% \def\ZSFDensityBoxes{1}         % Innen-/Aussenabstände der Inhaltsboxen
% \def\ZSFDensityText{1}          % Absatzabstand im Fliesstext
% \def\ZSFDensityTables{1}        % Zell- und Zeilenabstand in Tabellen
% \def\ZSFDensityStructure{1}     % Polsterung und Abstände der Titelbalken
%
% Freiraum, der vor einem Titelbalken übrig sein muss, sonst bricht die
% Spalte davor um. Grösser = Balken stehen seltener knapp am Spaltenfuss.
% \def\ZSFBreakReserveFactor{1}   % 1 = Referenz
%
% Der Abstand NACH einem Titelbalken — der ganze Abstand bis zum ersten
% Block darunter, Fliesstext wie Box.
% \def\ZSFBarGapFactor{1}         % 1 = Referenz (3pt), kleiner = enger
%
% Schriftgrösse der Diagramm-Beschriftung — beide Stufen gemeinsam.
% \def\ZSFDiagramLabelScale{0.85} % 1 = Referenz, kleiner = kleinere Labels

\input{styles/00_packages}
\input{styles/10_math}

% ── Formelsatz: Summen- und Limes-Grenzen ────────────────────────────
% side = neben dem Zeichen (spart Höhe), stack = darüber/darunter,
% auto = LaTeX entscheidet nach Kontext.
% \SetZSFsumMode{auto}            % auto | side | stack
% \SetZSFlimMode{auto}            % auto | side | stack

% ============================================================
% 11_math_advanced.tex — OPTIONALES Math-Erweiterungsmodul
% ============================================================
% Opt-in: in preamble.tex einkommentieren für mathe-lastige
% Fächer (LinAlg, Analysis). Baut auf 10_math.tex auf und
% ergänzt nur das Delta — keine Basis-Makros duplizieren.
%
% Voraussetzungen (aus 00_packages.tex):
%   - mathtools, unicode-math   (Operatoren, Symbole)
%   - tikz                      (drawbrace / tikzmark / annote)
% ============================================================

% --- Lineare-Algebra-Operatoren ---
\DeclareMathOperator{\Ker}{Ker}
\DeclareMathOperator{\rang}{rang}
\DeclareMathOperator{\Spur}{Spur}
\DeclareMathOperator{\diag}{diag}
\DeclareMathOperator{\spanop}{span}
\DeclareMathOperator{\eig}{eig}
\DeclareMathOperator{\algVfh}{algVfh}
\DeclareMathOperator{\gVfh}{gVfh}
\DeclareMathOperator{\proj}{proj}
\DeclareMathOperator{\sig}{sig}
\DeclareMathOperator{\Adj}{Adj}

% --- Gotisches Im/Re -> moderne aufrechte Operatoren ---
% AtBeginDocument, weil unicode-math seine Symbole ebenfalls erst dort
% definiert — eine Praeambel-Renewcommand wuerde still ueberschrieben.
\AtBeginDocument{%
  \let\oldIm\Im
  \renewcommand{\Im}{\operatorname{Im}}%
  \let\oldRe\Re
  \renewcommand{\Re}{\operatorname{Re}}%
}

% --- TikZ-Klammern für Matrix-Annotationen ---
\usetikzlibrary{decorations.pathreplacing,calc}
% Das optionale Argument sind KNOTEN-OPTIONEN (z.B. xshift), keine Länge. Der
% Default ist deshalb leer: Ein Mass an dieser Stelle liest pgf als
% Pfeilspezifikation und bricht ab ("Unknown arrow tip kind"). Der vertikale
% Versatz der Klammer hängt ohnehin an den Koordinaten, siehe unten.
\newcommand{\tikzmark}[2][]{\tikz[remember picture, overlay, baseline=-0.5ex]\node[#1](#2){};}
% Die Marker von \tikzmark sitzen auf der GRUNDLINIE des Terms. Ohne Versatz
% liegt die Klammer-Sehne genau dort und läuft als waagrechte Linie mitten
% durch die Formel. Der Versatz muss an den Koordinaten hängen, nicht als
% yshift an der Pfadoption: Knoten aus einem anderen Bild (remember picture)
% werden absolut aufgelöst und ignorieren die Pfad-Transformation.
%
% Richtung und Dekoration gehören zusammen, deshalb setzt jeder Stil beides.
% \ZSF@braceSign trägt das Vorzeichen zur Koordinate — die Optionen werden vor
% den Koordinaten ausgewertet, der Wert steht also rechtzeitig.
\makeatletter
\newcommand{\ZSF@braceSign}{-}
\tikzset{
  brace/.style={decorate, decoration={brace},
                /utils/exec=\renewcommand{\ZSF@braceSign}{}},
  brace mirrored/.style={decorate, decoration={brace,mirror},
                /utils/exec=\renewcommand{\ZSF@braceSign}{-}},
}
\newcounter{brace}
\setcounter{brace}{0}
% \drawbrace[stil]{marke links}{marke rechts}
%   brace           — Klammer ÜBER dem Term (Default)
%   brace mirrored  — Klammer UNTER dem Term
\newcommand{\drawbrace}[3][brace]{%
  \refstepcounter{brace}%
  \tikz[remember picture, overlay]\draw[#1]
    ([yshift=\ZSF@braceSign\ZSFbraceDrop]#2.center)--%
    ([yshift=\ZSF@braceSign\ZSFbraceDrop]#3.center)%
    node[pos=0.5, name=brace-\thebrace]{};%
}
\makeatother
\newcommand{\annote}[3][]{%
  \tikz[remember picture, overlay]\node[#1] at (#2) {#3};%
}

%
%
%
%
%
%
%
%
%
%
\newcommand{\ZSFbraceRoom}{\vspace{\dimexpr\ZSFbraceDrop+\ZSFspaceM\relax}}

% --- Asterisk-Definitionen (mathabx-Namen auf Standard-Symbole) ---
\newcommand{\asterisk}{\ast}
\newcommand{\oasterisk}{\circledast}

%
%
%
%
%
%
%
%
%
%
%
%
%
%
%
%
%
%
%
%
%
%
%
%
\makeatletter
\newcommand{\ZSF@augStretch}{\ZSFtableArrayStretch}
\newcommand{\ZSF@augOpen}{(}   % chktex 9
\newcommand{\ZSF@augClose}{)}  % chktex 9
\pgfkeys{
  /zsf/augmatrix/.is family, /zsf/augmatrix,
  rows/.is choice,
  rows/normal/.code={\renewcommand{\ZSF@augStretch}{\ZSFtableArrayStretch}},
  rows/roomy/.code ={\renewcommand{\ZSF@augStretch}{\ZSFtableArrayStretchRoomy}},
  rows/tight/.code ={\renewcommand{\ZSF@augStretch}{\ZSFtableArrayStretchTight}},
  delims/.is choice,
  delims/paren/.code  ={\renewcommand{\ZSF@augOpen}{(}\renewcommand{\ZSF@augClose}{)}},   % chktex 9
  delims/bracket/.code={\renewcommand{\ZSF@augOpen}{[}\renewcommand{\ZSF@augClose}{]}},   % chktex 9
  delims/none/.code   ={\renewcommand{\ZSF@augOpen}{.}\renewcommand{\ZSF@augClose}{.}},
  ZSF@augDefaults/.style={rows=normal, delims=paren},
}
\NewDocumentEnvironment{ZSFaugmatrix}{O{} m m}
  {\pgfkeys{/zsf/augmatrix, ZSF@augDefaults, #1}%
   \renewcommand{\arraystretch}{\ZSF@augStretch}%
   \left\ZSF@augOpen\begin{array}{@{}*{#2}{c}|*{#3}{c}@{}}}
  {\end{array}\right\ZSF@augClose}
\makeatother
   % opt-in: LinAlg — Ker/rang/Spur/proj/Im/Re,
%
%
%
%
%
%
%
%
%
\ifdefined\ZSFShowcaseMode
  \input{styles/12_plots}
\fi
% ============================================================
% 20_tables.tex — Tabellen-Helfer (semantisches Table-System)
% ============================================================
% Zentrales Design-Prinzip:
%   1. Tabellen fuellen IMMER \linewidth (kein manuelles p{0.16..}).
%   2. Spaltenbreiten werden PROPORTIONAL deklariert (Y{0.18} L{0.54} ...),
%      sodass die Summe der Faktoren = Anzahl X-Spalten ergibt.
%   3. Zebra-Streifen + Header-Highlight sind ins Environment eingebaut;
%      direkte Farb-Primitive sind in Kapiteln per Linter verboten.
%
% Damit ist strukturell ausgeschlossen:
%   - zu kurze Zebra-Streifen (Tabelle schmaler als Box),
%   - überstehende Zeilen (Content sprengt Box-Breite),
%   - Drift zwischen Kapiteln (jede Tabelle nutzt dieselben Wrapper).

% ------------------------------------------------------------
% Spaltentypen (tabularx-X-Varianten)
% ------------------------------------------------------------
% Feste Varianten — jede zaehlt als Faktor 1:
% m{#1}: X-Spalten vertikal mittig (Text + Grafik in einer Zeile).
\renewcommand{\tabularxcolumn}[1]{m{#1}}
\newcolumntype{L}{>{\RaggedRight\arraybackslash}X}   % linksbuendig, umbricht
\newcolumntype{C}{>{\Centering\arraybackslash}X}     % zentriert,  umbricht
\newcolumntype{R}{>{\RaggedLeft\arraybackslash}X}    % rechtsbuendig, umbricht

% Parametrisierte, proportionale Varianten — #1 = relativer Anteil.
%
% WICHTIG (tabularx-Konvention): Die Summe aller Faktoren in EINER
% Tabelle muss = Anzahl X-Spalten sein, damit die Tabelle \linewidth
% voll ausfuellt. L/C/R zählen als Faktor 1.
%
% Beispiele (2 Spalten → Summe 2):
%   Y{0.6}  Y{1.4}            → 30% / 70%
%   Y{1}    Y{1}              → 50% / 50% (== L L)
%   Q{0.8}  Y{1.2}            → 40% / 60%, Spalte 1 rechtsbuendig
% Beispiele (3 Spalten → Summe 3):
%   Y{0.6}  Y{1.5}  Y{0.9}    → 20% / 50% / 30%
%   L       Y{2}    L         → 25% / 50% / 25%
% Beispiele (4 Spalten → Summe 4):
%   Y{1.2}  Y{0.8}  Y{1.2}  Y{0.8}  → 30% / 20% / 30% / 20%
\newcolumntype{Y}[1]{>{\hsize=#1\hsize\RaggedRight\arraybackslash}X} % linksb., proport.
\newcolumntype{Z}[1]{>{\hsize=#1\hsize\Centering\arraybackslash}X}   % zentriert, proport.
\newcolumntype{Q}[1]{>{\hsize=#1\hsize\RaggedLeft\arraybackslash}X}  % rechtsb., proport.

% Zentrierte Formel-Spalte (Math-Mode automatisch, proportional)
\newcolumntype{F}[1]{>{\hsize=#1\hsize\Centering\arraybackslash$}X<{$}}

% ------------------------------------------------------------
% Luft zwischen Zellinhalt und Zeilenkante
% ------------------------------------------------------------
% Der Regler `rows` versprach »mehr Höhe für Zellen mit Brüchen« und konnte
% das nicht halten: \arraystretch skaliert das Zeilenstreben, und ein Streben
% ist ein MINDESTmass. Inhalt, der höher ist als es — ein \dfrac, eine Wurzel,
% ein Bild —, bleibt davon unberührt. Gemessen: an einer Textabelle +22 %
% Höhe, an derselben Tabelle mit \dfrac-Zellen +5 %.
%
% Bündig steht der Inhalt zusätzlich, und zwar unten. Grund ist die
% Ausrichtung der m-Spalte in array.sty (\ar@align@mcell): Ist eine Zelle
% höher als ein Streben, senkt sie den Block um
%   0.5 * (Zellhöhe - Strebenhöhe + \baselineskip)
% und schiebt ihn damit so weit nach unten, dass die Tiefe der Zelle die Tiefe
% der Zeile bestimmt — die ganze verbleibende Luft landet oben. Die Formel ist
% für mehrzeiligen Text gedacht; für eine einzelne hohe Zeile bedeutet sie
% Unterkante bündig. Im Satz sichtbar: Der Nenner eines \dfrac berührte die
% Zeilengrenze, in der letzten Zeile lief er aus der Box.
%
% Kein strebenbasierter Weg behebt das, weil jede Strebe ein Mindestmass ist
% und die Zeilentiefe das Maximum aus Strebe und Zellinhalt. Deshalb meldet
% der Ausrichter die Zelle um \ZSFtableCellPadY höher UND tiefer, ohne sie zu
% verschieben — dieselbe \raisebox-Geste wie bei \ZSFimage, nur eine Ebene
% höher und damit für JEDEN hohen Zellinhalt.
%
% Nur der hohe Zweig wird angefasst. Kurzer Text läuft weiter über die
% Grundlinien-Ausrichtung von array.sty; an ihm ändert sich nichts.
%
% \ar@align@mcell ist ein Hook und keine Interna-Bastelei: array.sty prüft ihn
% selbst mit \ifx\ar@align@mcell\@undefined und setzt ihn in anderen Pfaden
% auf \relax. Die Wache darunter hält den Fall offen, dass eine ältere
% array.sty ihn gar nicht kennt — dann bleibt alles beim Alten.
%
% Bewusst mit TeX-Primitiven und einem EIGENEN Register: \raisebox und die
% \@tempdim*-Register sind hier belegt — colortbl haelt den Farbueberhang der
% Zeile in \@tempdimb/\@tempdimc, und eine Zuweisung darauf wuerde die
% Zeilenflaeche verschieben. Mit \raisebox gemessen: die Zellen kamen leer
% heraus.
\newlength{\ZSFtableCellPadY}
\makeatletter
\newdimen\ZSF@cellShift
\@ifundefined{ar@align@mcell}{}{%
  \def\ar@align@mcell{%
    \ifdim \ht\ar@mcellbox > \ht\strutbox
      \ZSF@cellShift\ht\ar@mcellbox
      \advance\ZSF@cellShift-\ht\@arstrutbox
      \advance\ZSF@cellShift\baselineskip
      \ZSF@cellShift=.5\ZSF@cellShift
      \setbox\ar@mcellbox\hbox{\lower\ZSF@cellShift\box\ar@mcellbox}%
      \ht\ar@mcellbox\dimexpr\ht\ar@mcellbox+\ZSFtableCellPadY\relax
      \dp\ar@mcellbox\dimexpr\dp\ar@mcellbox+\ZSFtableCellPadY\relax
      \box\ar@mcellbox
    \else
      \box\ar@mcellbox
    \fi}%
}
\makeatother

% ------------------------------------------------------------
% Zebra-Streifen (Scannbarkeit)
% ------------------------------------------------------------
\newcommand{\ZSFzebraBG}{\TableZebraColor}
\newcommand{\SetZSFzebraBG}[1]{\renewcommand{\ZSFzebraBG}{#1}}
%
%
%
\newcommand{\ZSFzebraStart}[1]{\rowcolors{#1}{\BoxPlainBackColor}{\ZSFzebraBG}}
\newcommand{\ZSFzebraOff}{\rowcolors{1}{}{}}

% ------------------------------------------------------------
% Header-Highlight
% ------------------------------------------------------------
% Konvention:
%   \ZSFheaderRow   — startet die Header-Zeile und setzt den Zeilen-Hintergrund.
%   \ZSFhead{<Text>} — umschliesst den Inhalt JEDER Header-Zelle (Farbe + bold).
%
% Warum zweigeteilt? In tabularx-X-Spalten propagiert \color, das direkt nach
% \rowcolor am Zeilenanfang gesetzt wird, NICHT zuverlässig über die
% &-Grenzen. Die erste Zelle verliert die Farbe, andere behalten sie — je
% nach Spaltenbreite und Content (siehe SI-Praefix-Bug). Per-Zelle-Wrapper
% via \ZSFhead{...} ist robust und konsistent.
\newcommand{\ZSFheaderRowBG}{\chaptercolor}
\newcommand{\ZSFheaderRowFG}{\chaptercolortext}
\newcommand{\ZSFheaderRow}{\rowcolor{\ZSFheaderRowBG}}
% \ZSFInkOwned steht direkt neben dem \color, das den Kontrast setzt: Die
% Kopfzelle ist eine gesaettigte Flaeche mit eigener Textfarbe, eine
% Inline-Farbe darin waere unlesbar (40_colors_structure -> Ink-Vertrag).
\newcommand{\ZSFhead}[1]{{\color{\ZSFheaderRowFG}\ZSFInkOwned\ZSFfontTableHead #1}}

%
%
%
%
%
%
%
%
%
%
%
%
%
%
%
%
%
%
%
%
%
%
%
%
%
%
%
%
%
%
%
%
%
%
%
%
%
%
%
%
%
%
%
%
\makeatletter
\newcommand{\ZSF@spanBadAlign}[1]{%
  \PackageError{ZSF}{Unbekannte Ausrichtung '#1' in \string\ZSFspan}%
    {Gueltig sind L, C und R -- dieselben Buchstaben wie bei den
     Spaltentypen L/C/R.}%
}
\newcommand{\ZSFspan}[3]{%
  \if L#1\multicolumn{#2}{l}{#3}%
  \else\if R#1\multicolumn{#2}{r}{#3}%
  \else\if C#1\multicolumn{#2}{c}{#3}%
  \else\multicolumn{#2}{c}{\ZSF@spanBadAlign{#1}#3}%
  \fi\fi\fi
}
\makeatother

%
%
%
%
%
%
%
%
%
%
%
%
%
%
%
%
%
%
%
%
%
%
%
%
%
%
%
%
%
%
%
%
%
%
%
%
%
%
\makeatletter
\newif\ifZSF@tblHeader
\newif\ifZSF@tblZebra
\newcommand{\ZSF@tblFont}{\ZSFfontContent}
\newcommand{\ZSF@tblStretch}{\ZSFtableArrayStretch}
\newlength{\ZSF@tblColSep}

%
%
%
%
%
%
%
%
%
%
%
%
%
%
%
%
%
%
%
\newcommand{\ZSF@setRows}[1]{%
  \IfStrEqCase{#1}{%
    {normal}{\renewcommand{\ZSF@tblStretch}{\ZSFtableArrayStretch}%
             \setlength{\ZSFtableCellPadY}{\ZSFDensityTables\ZSFspaceXS}}%
    {roomy} {\renewcommand{\ZSF@tblStretch}{\ZSFtableArrayStretchRoomy}%
             \setlength{\ZSFtableCellPadY}{\ZSFDensityTables\ZSFspaceS}}%
    {tight} {\renewcommand{\ZSF@tblStretch}{\ZSFtableArrayStretchTight}%
             \setlength{\ZSFtableCellPadY}{0pt}}%
  }%
}
\newcommand{\ZSF@setColsep}[1]{%
  \IfStrEqCase{#1}{%
    {normal}{\setlength{\ZSF@tblColSep}{\ZSFtableEdgePad}}%
    {tight} {\setlength{\ZSF@tblColSep}{\ZSFtableEdgePadTight}}%
  }%
}
\newcommand{\ZSF@tblSpec}{}
\newcommand{\ZSF@beginTabularx}[1]{\tabularx{\linewidth}{#1}}
\pgfkeys{
  /zsf/table/.is family, /zsf/table,
  header/.is if=ZSF@tblHeader,
  zebra/.is if=ZSF@tblZebra,
  font/.is choice,
  font/normal/.code={\renewcommand{\ZSF@tblFont}{\ZSFfontContent}},
  font/dense/.code={\renewcommand{\ZSF@tblFont}{\ZSFfontContentDense}},
  % `rows` und `colsep` beantworten ihre Frage EINMAL — die beiden Schlüssel
  % hier reichen sie nur durch (\ZSF@setRows / \ZSF@setColsep, oben). Das
  % valuegrid stellt dieselben zwei Fragen und liest dieselbe Antwort; stünde
  % die Zuordnung an beiden Bausteinen, hiesse `rows=roomy` beim nächsten
  % Eingriff auf einem der beiden etwas anderes.
  rows/.is choice,
  rows/normal/.code={\ZSF@setRows{normal}},
  rows/roomy/.code ={\ZSF@setRows{roomy}},
  rows/tight/.code ={\ZSF@setRows{tight}},
  colsep/.is choice,
  colsep/normal/.code={\ZSF@setColsep{normal}},
  colsep/tight/.code ={\ZSF@setColsep{tight}},
  ZSF@tblDefaults/.style={header=true, zebra=true, font=normal,
                          rows=normal, colsep=normal},
}
\NewDocumentEnvironment{ZSFtable}{O{} m}{%
  \begingroup
  % Zellinhalt: ein \ZSFimage hier bekommt Zellen-, nicht Abbildungshoehe.
  % Explizit und nicht als Default verlassen, damit eine Tabelle IN einer
  % figbox nicht deren Abbildungsbudget erbt (siehe \ZSF@imgBudget, 60_boxes).
  \renewcommand{\ZSF@imgBudget}{\ZSFimageMaxHeight}%
  \pgfkeys{/zsf/table, ZSF@tblDefaults, #1}%
  \setlength{\tabcolsep}{\ZSF@tblColSep}%
  \renewcommand{\arraystretch}{\ZSF@tblStretch}%
  \ZSF@tblFont
  \ifZSF@tblZebra
    \ifZSF@tblHeader\ZSFzebraStart{2}\else\ZSFzebraStart{1}\fi
  \else
    \ZSFzebraOff
  \fi
%
%
%
%
%
%
%
%
%
%
%
%
%
%
%
%
%
%
%
%
%
%
%
%
%
%
%
%
%
%
%
%
%
%
%
  \edef\ZSF@tblSpec{@{\hskip\ZSF@tblColSep}#2@{\hskip\ZSF@tblColSep}}%
  \expandafter\ZSF@beginTabularx\expandafter{\ZSF@tblSpec}%
}{%
  \endtabularx
  \endgroup
}
\makeatother


% ── Tabellen: Zebra-Ton ──────────────────────────────────────────────
% Default folgt der Kapitelfarbe. Ein fester Ton macht Tabellen über alle
% Kapitel hinweg gleich — sinnvoll, wenn die Kapitelfarbe schon woanders trägt.
% \SetZSFzebraBG{\TableZebraColor}

% ============================================================
% 30_layout_spacing.tex — Layout-Grundlagen, Spacing-Skala, Box-Bindungen
% ============================================================

\setlength{\parindent}{0pt}

\widowpenalty=10000
\clubpenalty=10000
\brokenpenalty=4000
\predisplaypenalty=10000

\raggedcolumns
\newcounter{zsfFirstChap}\setcounter{zsfFirstChap}{1}

%
%
%
%
%
%
%
%
%
%
%
%
%
%
%
%
%
%
%
\providecommand{\ZSFDensityFactor}{1}

% ── \ZSFderive: eine Ableitung, die man wiederholen kann ─────────────
% Setzt das Register sofort UND merkt sich die Rechnung. \ZSFChapterDensity
% (unten) laesst die gemerkte Kette in derselben Reihenfolge noch einmal
% laufen und bekommt damit den ganzen Registersatz auf einen neuen Faktor.
%
% Der Grund fuer die Kette statt einer Handvoll Neuberechnungen: Ein Register
% ist eine LAENGE, kein Ausdruck. \ZSFboxPadX hat den Wert, den \ZSFspaceS bei
% seiner Berechnung hatte — die vier Basisregister nachtraeglich zu aendern
% erreicht es nicht mehr. Wer eine neue abgeleitete Groesse anlegt, benutzt
% deshalb \ZSFderive und nicht \setlength; sie faehrt dann von selbst mit.
\makeatletter
\newcommand{\ZSF@spacingChain}{}
\newcommand{\ZSFderive}[2]{%
  \setlength{#1}{#2}%
  \expandafter\gdef\expandafter\ZSF@spacingChain\expandafter{%
    \ZSF@spacingChain\setlength{#1}{#2}}%
}
% Der Dokumentwert, auf den \ZSFChapterDensityReset zurueckstellt. Muss VOR der
% ersten Ableitung gesichert werden, sonst ist er beim Zuruecksetzen schon weg.
\let\ZSF@densityDocument\ZSFDensityFactor

%
%
%
%
%
%
%
%
%
%
%
%
\newcommand{\ZSFChapterDensity}[1]{\def\ZSFDensityFactor{#1}\ZSF@spacingChain
  \ZSFApplyDisplaySkips}
\newcommand{\ZSFChapterDensityReset}{%
  \let\ZSFDensityFactor\ZSF@densityDocument \ZSF@spacingChain
  \ZSFApplyDisplaySkips}
\makeatother

% ── Display-Skips: die EINE Stelle, und sie muss wiederholbar sein ───
% Die engen Abstände um Formeln sind eine ZSF-Entscheidung. LaTeX legt sie
% aber in den Groessenbefehlen ab: \normalsize und \footnotesize setzen alle
% vier Skips auf die KLASSENWERTE zurueck. Jeder Schriftwechsel im Inhalt
% loescht damit die Entscheidung — lautlos, denn Formeln sehen auch mit den
% Klassenwerten korrekt aus, nur luftiger.
%
% Deshalb ist die Politik ein benanntes Makro und kein einmaliges \setlength:
% Wer die Groesse wechselt, ruft sie danach erneut auf. Genutzt von
% \ZSFfontContent/\ZSFfontContentDense (der `font`-Regler), vom
% \AtBeginDocument in 70_document_settings und von \ZSFChapterDensity oben.
% \jot und \arraycolsep laufen hier mit, obwohl kein Groessenbefehl sie
% zuruecksetzt: Sie gehoeren zur selben Frage — »wie viel Luft hat eine
% Formel« — und eine zweite Anwendungsstelle waere eine zweite Auslegung.
\newcommand{\ZSFApplyDisplaySkips}{%
  \setlength{\abovedisplayskip}{\ZSFspaceXS}%
  \setlength{\belowdisplayskip}{\ZSFspaceXS}%
  \setlength{\abovedisplayshortskip}{\ZSFspaceXS}%
  \setlength{\belowdisplayshortskip}{\ZSFspaceXS}%
  \setlength{\mathsurround}{\ZSFspaceXS}%
  \setlength{\jot}{\ZSFmathRowSep}%
  \setlength{\arraycolsep}{\ZSFmathColSep}%
}

\newlength{\ZSFspaceXS}\ZSFderive{\ZSFspaceXS}{\ZSFDensityFactor\dimexpr 1pt\relax}
\newlength{\ZSFspaceS} \ZSFderive{\ZSFspaceS}{\ZSFDensityFactor\dimexpr 4pt\relax}
\newlength{\ZSFspaceM} \ZSFderive{\ZSFspaceM}{\ZSFDensityFactor\dimexpr 7pt\relax}
\newlength{\ZSFspaceL} \ZSFderive{\ZSFspaceL}{\ZSFDensityFactor\dimexpr 10pt\relax}

%
%
%
%
%
%
%
%
%
%
%
%
%
%
%
%
%
%
%
\providecommand{\ZSFDensityBoxes}{1}
\providecommand{\ZSFDensityText}{1}
\providecommand{\ZSFDensityTables}{1}
\providecommand{\ZSFDensityStructure}{1}

% Outer (Text-Box-Bindung) und Sep (Trennlinien) — beide folgen der Dichte.
% Outer ist bewusst als \ZSFspaceM abgeleitet statt erneut ausgeschrieben:
% es IST der Standard-Aussenabstand, kein eigener Wert.
\newlength{\ZSFspaceOuter}\ZSFderive{\ZSFspaceOuter}{\ZSFspaceM}
\newlength{\ZSFspaceSep}  \ZSFderive{\ZSFspaceSep}{\ZSFDensityBoxes\dimexpr
  \ZSFDensityFactor\dimexpr 3pt\relax\relax}

% Absatzabstand im Fliesstext — der grösste einzelne vertikale Posten einer
% ZSF und deshalb der wichtigste Teil der Dichte-Stellschraube. Muss nach den
% Basisregistern stehen, damit \ZSFDensityFactor bereits bekannt ist.
% Der Bereichsfaktor greift über das Zwischenregister: Ein Faktor vor einem
% Skip-Register skaliert Soll-, Streck- und Schrumpfanteil gemeinsam; zwei
% Faktoren direkt hintereinander kann TeX dagegen nicht lesen.
\newlength{\ZSFparskipBase}
\ZSFderive{\ZSFparskipBase}{\ZSFDensityFactor\dimexpr 3.5pt\relax
                    plus \ZSFDensityFactor\dimexpr 1pt\relax
                   minus \ZSFDensityFactor\dimexpr 0.5pt\relax}
\ZSFderive{\parskip}{\ZSFDensityText\ZSFparskipBase}

% ── Zeilenhöhe: die zweite globale Stellschraube ─────────────────────
% Bewusst NICHT an \ZSFDensityFactor gekoppelt. Abstände sind Leerraum und
% vertragen jede Skalierung; Zeilenhöhe enthält die Schrift selbst — sie mit
% demselben Faktor zu stauchen, lässt Formeln und Akzente kollidieren.
% Wer zusätzlich zur Dichte enger setzen will, dreht hier separat:
%   \def\ZSFLeadingFactor{0.92}   % in preamble.tex, vor diesem \input
\providecommand{\ZSFLeadingFactor}{0.95}
\linespread{\ZSFLeadingFactor}

% ------------------------------------------------------------
% Spalten-Layout
% ------------------------------------------------------------
% \columnsep folgt der Dichte bewusst NICHT: Beim Verdichten sollen die
% Zeilen enger stehen, nicht die Spalten zusammenrücken — ein schmalerer
% Steg macht die ZSF schlechter lesbar, obwohl er kaum Platz spart.
\newlength{\ZSFcolumnSep}\setlength{\ZSFcolumnSep}{10pt}
\setlength{\columnsep}{\ZSFcolumnSep}
\setlength{\multicolsep}{\ZSFspaceS}
\setlength{\emergencystretch}{2em}

% ------------------------------------------------------------
% Schwellwerte und Titelbalken-Hilfsgrössen
% ------------------------------------------------------------
% Mindesthöhe, die vor einem Titelbalken frei sein muss — sonst bricht die
% Spalte davor um. Die Reserve muss den Balken UND den Anfang der folgenden
% Box tragen, sonst bleibt der Balken allein am Spaltenende stehen.
% Deshalb staffeln sich die Werte nach der Grösse der eröffneten Einheit:
% Kapitel > Abschnitt > Unterabschnitt.
% Die drei Werte sind EINE Staffelung, kein Trio unabhängiger Zahlen. Wer die
% ZSF früher oder später umbrechen lassen will, dreht deshalb den Faktor und
% nicht die drei Register einzeln — nur so bleibt ihr Verhältnis gewahrt.
% Grösser = die Spalte bricht eher vorzeitig um und lässt
% Balken seltener knapp am Fuss stehen; kleiner = dichter gefüllte Spalten.
\providecommand{\ZSFBreakReserveFactor}{1}
\newlength{\ZSFbreakThresholdChapter}
  \setlength{\ZSFbreakThresholdChapter}{\ZSFBreakReserveFactor\dimexpr 10\baselineskip\relax}
\newlength{\ZSFbreakThreshold}
  \setlength{\ZSFbreakThreshold}{\ZSFBreakReserveFactor\dimexpr 7\baselineskip\relax}
\newlength{\ZSFbreakThresholdSubsubsection}
  \setlength{\ZSFbreakThresholdSubsubsection}{\ZSFBreakReserveFactor\dimexpr 5\baselineskip\relax}
\newlength{\ZSFtitleBarHeight}\setlength{\ZSFtitleBarHeight}{9.2pt}
\newlength{\ZSFsubsectionBarBeforeSkip}\ZSFderive{\ZSFsubsectionBarBeforeSkip}{\ZSFDensityStructure\ZSFspaceL}
\newlength{\ZSFsubsubsectionBarBeforeSkip}\ZSFderive{\ZSFsubsubsectionBarBeforeSkip}{\ZSFDensityStructure\ZSFspaceS}
\newlength{\ZSFsubsectionAfterChapterGap}\ZSFderive{\ZSFsubsectionAfterChapterGap}{\ZSFDensityStructure\ZSFspaceXS}

% ── Der Abstand NACH einem Titelbalken ───────────────────────────────
% EIN Register für alle vier Balken, und es ist der GANZE Abstand von der
% Unterkante des Balkens bis zur Oberkante des folgenden Blocks — Fliesstext
% wie Box. Der Weg dorthin ist \ZSFInterlude (unten), aufgerufen im
% after-Hook \ZSFTitleBarAfter (60_boxes).
%
% Vorher war dieser Abstand herrenlos, und zwar nach demselben Muster wie
% \ZSFBoxBeforeSkip (styles/README → »Abstand hat genau einen Schreiber«):
% Zwei Register (\ZSFchapterBarAfterSkip 2pt, \ZSFsubsectionAfterGap 1pt)
% setzten je EINEN Posten von dreien. Dazu kamen \parskip (3.5pt, weil der
% Text nach dem Balken ein neuer Absatz ist) und die Zeilenschaltung
% (\baselineskip − Höhe der ersten Zeile). Gemessen am Satz: Balken → Box
% 1pt, Balken → Fliesstext 6.8pt, Kapitelbalken → Fliesstext 9.5pt. Drei
% Abstände an derselben Stelle des Systems, keiner davon gewählt.
%
% Der Wert (3pt) liegt zwischen den beiden Zuständen von vorher: enger als der
% Fliesstext-Fall (6.8/9.5pt, den niemand gewählt hatte), etwas luftiger als
% der Box-Fall (1pt, der nur deshalb so knapp war, weil dort kein \parskip
% dazukam). Am Satz geprüft, nicht gerechnet: 2pt drückt den ersten Satz an
% den Balken, 4pt bringt die alte Lücke zurück.
%
% Die Staffelung Kapitel > Abschnitt ist bewusst aufgegeben: Die Ebene wird
% über Balkenhöhe, Farbe und den Abstand DAVOR unterschieden
% (\ZSFsubsectionBarBeforeSkip), nicht über einen Punkt danach.
\providecommand{\ZSFBarGapFactor}{1}
\newlength{\ZSFbarAfterGap}\ZSFderive{\ZSFbarAfterGap}{\ZSFBarGapFactor\dimexpr
  \ZSFDensityStructure\dimexpr\ZSFDensityFactor\dimexpr 3pt\relax\relax\relax}
% Der sichtbare Aussenradius ist Form, nicht Leerraum: Auch bei dichterem Satz
% bleiben die Kanten bewusst gleich weich. Der Innenradius wird aus der
% jeweiligen Rahmenbreite abgeleitet, damit die Kontur in der Kurve nicht
% dicker wird als auf den geraden Kanten.
\newlength{\ZSFboxArc}\setlength{\ZSFboxArc}{3pt}
\newlength{\ZSFboxRuleWidth}\setlength{\ZSFboxRuleWidth}{0.5pt}
\newlength{\ZSFboxInnerArc}\ZSFderive{\ZSFboxInnerArc}{\dimexpr
  \ZSFboxArc-\ZSFboxRuleWidth\relax}
\newlength{\ZSFgoalConditionRuleWidth}\setlength{\ZSFgoalConditionRuleWidth}{1.5pt}
\newlength{\ZSFgoalConditionBorderlineWidth}\setlength{\ZSFgoalConditionBorderlineWidth}{0.2pt}
\newlength{\ZSFgoalConditionInnerArc}\ZSFderive{\ZSFgoalConditionInnerArc}{\dimexpr
  \ZSFboxArc-\ZSFgoalConditionRuleWidth\relax}
% Abstand zwischen zwei Gliedern einer Zielkette (Bedingung, Schritt, Ziel).
% Ein eigenes Register, weil der Abstand in der goalbox eine ENTSCHEIDUNG ist
% und nicht das Nebenprodukt der Zeilenschaltung: Die Glieder sind verschieden
% hoch (Pille, Displayzeile, gerahmtes Ziel), und ohne einen benannten Wert
% ergibt sich der Rhythmus aus der Tinte statt aus dem Satz (siehe
% \ZSFInterlude). Gleich \ZSFboxPadY, damit der Abstand zum Rahmen derselbe
% ist wie der zwischen den Gliedern — die Kette liest sich dadurch als ein
% gleichmässiger Stapel.
\newlength{\ZSFgoalLinkSep}  \ZSFderive{\ZSFgoalLinkSep}{\ZSFDensityBoxes\ZSFspaceS}
% Zeilenhöhe der ZSFtable in drei Stufen — der Regler `rows` wählt zwischen
% ihnen (20_tables). Drei Werte statt eines, weil die Höhe einer Zeile nicht
% aus der Tabellenart folgt, sondern aus ihrem Inhalt: Eine Tabelle mit
% \dfrac-Zellen braucht mehr Höhe, ein Symbolregister mit Einzelzeichen
% weniger. `font=dense` ist dafür die falsche Antwort — es verkleinert die
% Schrift, während der Bedarf hier die Höhe ist.
%
% Als Faktoren (\newcommand) statt als Längen: \arraystretch ist ein
% Multiplikator auf die Zeilenhöhe, kein Mass. Deshalb auch bewusst NICHT
% über \ZSFDensityTables skaliert — der Bereichsfaktor wirkt auf die
% Abstände der Tabelle (\ZSFtableEdgePad, \extrarowheight), und ein Faktor
% mal einem Faktor ergäbe eine Stauchung, die niemand vorhersagt.
\newcommand{\ZSFtableArrayStretch}{1.04}
\newcommand{\ZSFtableArrayStretchRoomy}{1.28}
\newcommand{\ZSFtableArrayStretchTight}{0.92}

% Maximale Höhe für Bilder in figboxen (verhindert überhohe Aspect-Ratios).
\newlength{\ZSFfigMaxHeight}\setlength{\ZSFfigMaxHeight}{2.6cm}

% Höhe eines Bildes OHNE Box (\ZSFimage) in einer TABELLENZELLE. Eigener Wert
% und nicht \ZSFfigMaxHeight, weil die beiden verschiedene Grössenordnungen
% sind: eine Abbildung füllt einen Block, ein Zellenbild muss in eine
% Tabellenzeile passen.
%
% Der eigentliche Zweck ist die Gleichmässigkeit: Eine Spalte aus
% \ZSFimage{pfad}-Aufrufen ohne Höhenangabe ist damit von selbst einheitlich
% hoch. Stünde die Höhe stattdessen an jedem Aufruf, wäre dieselbe Entscheidung
% fünfmal geschrieben — und beim sechsten Bild einmal anders.
%
% Welcher der beiden Werte gilt, entscheidet NICHT das Bild, sondern der
% Baustein, in dem es steht: ZSFtable und valuegrid binden auf diesen Wert,
% figbox und splitbox auf \ZSFfigMaxHeight (Mechanik: \ZSF@imgBudget in
% styles/60_boxes.tex). Ein fester Default hier wäre an einem der beiden Orte
% still falsch — und still ist das Wort, das zählt: Ein zu klein gesetztes Bild
% bricht keinen Build und löst keine Warnung aus.
\newlength{\ZSFimageMaxHeight}\setlength{\ZSFimageMaxHeight}{1.1cm}

% Es gab hier ein \ZSFimagePadY — die senkrechte Luft, die ein Bild in einer
% Tabellenzelle braucht, weil ein \includegraphics Höhe und keine Tiefe hat
% und dadurch an der Oberkante der Zeile klebt. Der Wert ist ENTFALLEN: Es ist
% derselbe Satz wie »der Inhalt einer Zeile hält Abstand zur Zeilenkante«, und
% der gilt für einen \dfrac genauso. Zwei Register für eine Aussage sind eine
% Kopie (styles/README → Vereinigungstest); geblieben ist \ZSFtableCellPadY
% in 20_tables, gesetzt vom Regler `rows`.

% Nutzbarer Anteil der Zeilenbreite in mehrteiligen Bild-Layouts; der Rest ist
% der \hfill-Steg zwischen den Elementen (Bildreihe in \ZSFfig, Bild-neben-Text
% in \ZSFfigside) — EIN Wert, damit die beiden Layouts nicht um einen Punkt
% auseinanderdriften. Als Makro statt \newlength: es ist ein Faktor, kein
% Mass, und wird erst an der Verwendungsstelle gegen die dortige \linewidth
% ausgewertet.
\newcommand{\ZSFfigUsableFraction}{0.97}

% Abstand der TikZ-Spannklammer (\drawbrace, Opt-in 11_math_advanced) von der
% Grundlinie des annotierten Terms. Ohne ihn läuft die Klammer als waagrechte
% Linie mitten durch die Formel, weil die Marker auf der Grundlinie sitzen.
\newlength{\ZSFbraceDrop}\ZSFderive{\ZSFbraceDrop}{\ZSFspaceM}

%
%
%
%
%
%
%
%
%
%
%
%
%
%
%
\newlength{\ZSFmathRowSep}\ZSFderive{\ZSFmathRowSep}{\ZSFspaceSep}
\newlength{\ZSFmathColSep}\ZSFderive{\ZSFmathColSep}{\ZSFDensityText\ZSFspaceS}

% ── Der Abstand von der Spaltenoberkante zur ersten Grundlinie ───────
% \topskip beim Anfang einer Spalte, \splittopskip am Anfang einer
% fortgesetzten. Beide standen auf den Klassenvorgaben (9pt und 10pt) und
% waren damit die einzigen Masse, die JEDE Spalte anfasst und die niemand
% gesetzt hat — sie sind nicht »klein«, sie waren nur unsichtbar, weil eine
% Spalte meist mit einem Balken beginnt und der die Glue aufbraucht.
%
% EIN Register für beide: Es ist dieselbe Frage, nur einmal am Anfang und
% einmal in der Fortsetzung. Und bewusst an \baselineskip statt an der
% Dichte — der Wert muss die Oberlänge der ersten Zeile tragen, sonst rutscht
% sie nach oben aus dem Satzspiegel. Damit folgt er \ZSFLeadingFactor und
% nicht \ZSFDensityFactor, aus demselben Grund wie die Zeilenhöhe selbst.
% Gesetzt wird er in 70_document_settings, wo \baselineskip feststeht.
\newlength{\ZSFcolumnTopSkip}

% Tokens für die ruhige Box-Familie und ZSFlist.
% Der Einzug ist der einzige Wert, der sich zwischen nummerierter und
% unnummerierter Liste unterscheidet — eine Zahl mit Punkt braucht mehr Platz
% als ein Aufzählungszeichen. Deshalb zwei Werte unter EINEM Namen statt
% zweier Registersätze mit verschiedenen Präfixen.
\newlength{\ZSFstatementBarWidth}\setlength{\ZSFstatementBarWidth}{2pt}
\newlength{\ZSFlistMarkGap}         \ZSFderive{\ZSFlistMarkGap}{\ZSFspaceXS}
\newlength{\ZSFlistItemSep}         \ZSFderive{\ZSFlistItemSep}{\ZSFspaceXS}
\newlength{\ZSFlistLeftMargin}      \setlength{\ZSFlistLeftMargin}{1.2em}
\newlength{\ZSFlistLeftMarginOrdered}\setlength{\ZSFlistLeftMarginOrdered}{1.6em}
\newlength{\ZSFlistLabelSep}        \setlength{\ZSFlistLabelSep}{0.4em}

% Kleinstes vertikales Innenmass für Text an einer Kante: Zeilenluft in
% Tabellen (\extrarowheight in 70_document_settings) und die knappe Polsterung
% jeder Box mit `pady=tight` (leise Box, goalbox, valuegrid).
\newlength{\ZSFtextMinPad}\ZSFderive{\ZSFtextMinPad}{\ZSFDensityBoxes\ZSFspaceXS}

%
%
%
%
%
%
%
%
%
%
%
%
%
%
%
%
%
%
%
%
%
%
%
%
%
\newlength{\ZSFboxPadX}      \ZSFderive{\ZSFboxPadX}{\ZSFDensityBoxes\ZSFspaceS}
\newlength{\ZSFboxPadXBar}   \ZSFderive{\ZSFboxPadXBar}{\ZSFDensityBoxes\ZSFspaceM}
\newlength{\ZSFtableEdgePad} \ZSFderive{\ZSFtableEdgePad}{\ZSFDensityTables\ZSFspaceM}
\newlength{\ZSFtableEdgePadTight}
  \ZSFderive{\ZSFtableEdgePadTight}{\ZSFDensityTables\ZSFspaceS}

% Vertikal dasselbe: Innenabstand oben/unten und der Abstand NACH einer Box.
% Beide als Register, damit der Abstand nach einer Box nicht pro Box-Stil
% erneut ausgeschrieben wird.
\newlength{\ZSFboxPadY}      \ZSFderive{\ZSFboxPadY}{\ZSFDensityBoxes\ZSFspaceS}
\newlength{\ZSFBoxAfterSkip} \ZSFderive{\ZSFBoxAfterSkip}{\ZSFDensityBoxes\ZSFspaceS}

% Balken haben ihren eigenen vertikalen Rhythmus — bewusst getrennt vom
% Box-Innenabstand, damit sich Balkenhoehe und Boxpolsterung unabhaengig
% einstellen lassen.
\newlength{\ZSFbarPadY}      \ZSFderive{\ZSFbarPadY}{\ZSFDensityStructure\ZSFspaceS}
\newlength{\ZSFbarPadYTight} \ZSFderive{\ZSFbarPadYTight}{\ZSFDensityStructure\ZSFspaceXS}

% ------------------------------------------------------------
% Code-Längen — nur für die optionalen Code-Module benötigt
% (styles/65_code_style.tex + 67_code_comments.tex). Schadlos
% wenn die Module in preamble.tex deaktiviert sind.
% Threshold:  Schwelle für Auto-Overflow (Code+Kommentar > Threshold -> oben)
% Das Default-Register ist immutabel und dient \ResetCodeCommentThreshold
% als Anker; \SetCodeCommentThreshold modifiziert nur das "lebende" Register.
% ------------------------------------------------------------
% Zeilenhoehe im Code-Block. Eigenes Register statt einer nackten Zahl im
% Listings-Stil: Sie ist die einzige Zeilenhoehen-Entscheidung neben
% \ZSFLeadingFactor und soll wie dieser auffindbar und drehbar sein.
\providecommand{\ZSFcodeLeadingFactor}{0.8}
% \ZSFcodeMiniSkip ist ENTFALLEN. Es war der Abstand um ein Listing im
% globalen \lstset — der Stil CodeExpert setzte an derselben Stelle
% \ZSFspaceM, und welcher Wert galt, hing daran, ob der Autor den Stil
% angefordert hatte. Seit der Stil global gilt (65_code_style), gibt es nur
% noch eine Antwort, und sie kommt aus der Skala.
\newlength{\ZSFcodeFrameRule} \setlength{\ZSFcodeFrameRule}{0.4pt}
\newlength{\ZSFcodeXMargin}   \setlength{\ZSFcodeXMargin}{8pt}
\newlength{\ZSFcodeCommentThresholdDefault}
  \setlength{\ZSFcodeCommentThresholdDefault}{\maxdimen}
\newlength{\ZSFcodeCommentThreshold}\setlength{\ZSFcodeCommentThreshold}{\ZSFcodeCommentThresholdDefault}
\newlength{\ZSFcodeCommentGap}      \ZSFderive{\ZSFcodeCommentGap}{\ZSFspaceM}

% Valuegrid, Index, Readability und dokumentweite Overlays
% \ZSFvaluegridRowSep ist ENTFALLEN. Es war der Zeilenabstand des Wertegitters
% und damit ein zweites Register für dieselbe Frage, die der Regler `rows` an
% der ZSFtable beantwortet. Beide Tabellen-Bausteine lesen jetzt
% \ZSFtableCellPadY (20_tables) — ein Wert, zwei Engines.
\newlength{\ZSFreadabilityEmergencyStretch}\setlength{\ZSFreadabilityEmergencyStretch}{2.5em}
% Flatterzone im Fliesstext: schmaler als der ragged2e-Standard (0pt plus 2em),
% damit die schmalen 4 Spalten gut gefüllt bleiben. Wirkt nur über
% \ZSFReadableOn — Tabellen-Spalten (L/Y) und Index behalten den Standard.
\newlength{\ZSFreadabilityRaggedRightskip}\setlength{\ZSFreadabilityRaggedRightskip}{0pt plus 1em}
\newlength{\ZSFindexItemIndent}       \setlength{\ZSFindexItemIndent}{1.2em}
\newlength{\ZSFindexSubitemIndent}    \setlength{\ZSFindexSubitemIndent}{2em}
\newlength{\ZSFindexSubitemOffset}    \setlength{\ZSFindexSubitemOffset}{1em}
\newlength{\ZSFindexSubsubitemIndent} \setlength{\ZSFindexSubsubitemIndent}{2.6em}
\newlength{\ZSFindexSubsubitemOffset} \setlength{\ZSFindexSubsubitemOffset}{1.8em}
\newlength{\ZSFindexParStretch}       \setlength{\ZSFindexParStretch}{0.3pt}
% Deckkraft der Release-Kennung im Footer. Als Makro, weil \transparent einen
% einheitenlosen Faktor nimmt und ein solcher Wert sonst durch jede Pruefung faellt.
\providecommand{\ZSFfooterOpacity}{0.5}
\newlength{\ZSFfooterBottomInset}     \setlength{\ZSFfooterBottomInset}{1mm}
\newlength{\ZSFfooterRightInset}      \setlength{\ZSFfooterRightInset}{2mm}
% Freiraum am Satzspiegel-Fuss für das Footer-Overlay (Release-Kennung und
% Seitenzahl). Das Overlay zeichnet ausserhalb des Textflusses und reserviert
% deshalb KEINE Höhe — ohne diesen Freiraum überdruckt es die letzte Zeile
% einer voll gefüllten Spalte. Der Wert wird in 70_document_settings aus der
% Footer-Schrift gemessen und an geometry übergeben; hier steht nur das
% Register, damit alle Masse an einem Ort liegen.
\newlength{\ZSFfooterClearance}
\newlength{\ZSFitemizeLeftMargin}     \setlength{\ZSFitemizeLeftMargin}{3mm}
\newlength{\ZSFenumerateLeftMargin}   \setlength{\ZSFenumerateLeftMargin}{4mm}
\newlength{\ZSFidentityMarkerInset}   \setlength{\ZSFidentityMarkerInset}{2mm}

%
%
%
%
%
%
%
%
%
%
%
%
%
\newif\ifzsfAfterChapterBar
\zsfAfterChapterBarfalse
\newcommand{\ZSFMarkAfterChapterBar}{\global\zsfAfterChapterBartrue}
\newcommand{\ZSFConsumeAfterChapterBar}{\global\zsfAfterChapterBarfalse}

% Der Abstand VOR einem Block. Direkt nach einem Titelbalken der Balken-Abstand,
% sonst der normale Blockabstand.
%
% Gelesen wird die Balken-SPERRE (\ifzsfTitleBarPending, 60_boxes) und kein
% eigener Merker: Sie sagt bereits genau das, was hier gebraucht wird — »ein
% Balken steht da, und es kam noch kein Inhalt«. Ein zweiter Merker daneben
% wäre dieselbe Aussage ein zweites Mal, und die Sperre gilt für ALLE vier
% Balken, während der frühere \ifzsfAfterSubsectionBar nur die beiden
% Abschnittsebenen kannte — nach einem Kapitelbalken zog derselbe Block
% deshalb einen anderen Abstand.
\newcommand{\ZSFBoxBeforeSkip}{%
  \ifzsfTitleBarPending\ZSFbarAfterGap\else\ZSFspaceS\fi
}

% Kollabiert den before-skip der SubsectionBar auf \ZSFsubsectionAfterChapterGap,
% wenn sie direkt nach einer ChapterBar steht.
\newcommand{\ZSFCollapseAfterChapterBar}{%
  \ifzsfAfterChapterBar
    \vspace*{\dimexpr\ZSFsubsectionAfterChapterGap-\ZSFsubsectionBarBeforeSkip\relax}%
    \ZSFConsumeAfterChapterBar
  \fi
}

% ============================================================
% TEXT-BOX-BINDUNGEN
% Koppeln Fliesstext direkt an Boxen — verhindert visuellen Bruch.
% ============================================================

% Robustes Skip-/Penalty-Aufrollen (vmode-/hmode-aware)
\newcommand{\ZSFRobustUnskip}{%
  \ifvmode
    \ifdim\lastskip>0pt\relax\vskip-\lastskip\fi
    \ifnum\lastpenalty=0\else\unpenalty\fi
    \ifdim\lastskip>0pt\relax\vskip-\lastskip\fi
  \else
    \ifdim\lastskip>0pt\relax\unskip\fi
    \ifnum\lastpenalty=0\else\unpenalty\fi
    \ifdim\lastskip>0pt\relax\unskip\fi
  \fi
}

%
%
%
%
%
%
\makeatletter
% Hier deklariert, nicht in 60_boxes: Das Flag gehört zu diesen beiden Makros
% und muss vor ihnen existieren, egal in welcher Reihenfolge die Module
% wachsen. Gesetzt wird es von der formulabox.
\newif\ifZSF@inFormulaBox
%
%
%
%
%
%
%
%
%
%
%
%
%
\newcommand{\textVorBox}[1]{%
  \ifZSF@inFormulaBox
    \par
    \begingroup
    \setlength{\parskip}{0pt}%
    \noindent\raggedright\ZSFBoxNote{#1}\par
    \endgroup
    \vspace{\ZSFspaceXS}%
    \ZSFBoxAlign
  \else
    \par
    \addvspace{\ZSFspaceOuter}%
    \begingroup
    \setlength{\parskip}{0pt}%
    {\ZSFfontBoxBinding\color{TextVorBoxColor} \noindent #1\par}%
    \endgroup
    \nopagebreak
    \vspace{\ZSFspaceXS}%
    \vspace{-\ZSFspaceS}%
    \nopagebreak
  \fi
}

% Text der zur Box DARÜBER gehört (Ergänzung/Folgerung)
\newcommand{\textNachBox}[1]{%
  \ifZSF@inFormulaBox
    \par
    \vspace{\ZSFspaceXS}%
    \begingroup
    \setlength{\parskip}{0pt}%
    \noindent\raggedright\ZSFBoxNote{#1}\par
    \endgroup
    \ZSFBoxAlign
  \else
    \par
    \ZSFRobustUnskip
    \vspace{\ZSFspaceXS}%
    \begingroup
    \setlength{\parskip}{0pt}%
    \nopagebreak
    {\ZSFfontBoxBinding\color{TextNachBoxColor} \noindent #1\par}%
    \endgroup
    \addvspace{\ZSFspaceOuter}%
  \fi
}
\makeatother

%
%
%
%
%
%
\NewDocumentCommand{\ZSFgap}{O{M}}{\vspace{\csname ZSFspace#1\endcsname}}
% \ZSFSectionGap bleibt als eigener Name, weil er etwas anderes ausdrückt als
% eine Grösse: den Themenwechsel. Dass er aktuell derselbe Abstand wie
% \ZSFgap[L] ist, ist eine Einstellung dieser Zeile, keine Verdopplung.
\newcommand{\ZSFSectionGap}{\ZSFgap[L]}

% ------------------------------------------------------------
% Vertikaler Einschub INNERHALB einer Box
% ------------------------------------------------------------
% Formelnote, Trennlinie, Item-Heading, Folgerung und die Glieder einer
% Zielkette setzen alle denselben Einschub. Als Makro statt von Hand, damit der
% Abstand an der Dichte-Stellschraube hängt und der Rhythmus an einer Stelle
% änderbar bleibt.
%
% DER ABSTAND IST EIN OPTISCHER, KEIN LEIM. Das ist der ganze Punkt dieser drei
% Zeilen und der Grund, warum sie nicht `\vspace{#1}` heissen:
%
% Zwischen zwei Blöcken im vertikalen Satz liegt nicht nur die gesetzte Glue,
% sondern zusätzlich die Zeilenschaltung — TeX schiebt vor jede Box so viel
% Leim ein, dass der Grundlinienabstand \baselineskip ergibt. Der SICHTBARE
% Abstand ist damit `#1 + (\baselineskip - Höhe der folgenden Zeile)`. Solange
% alle Zeilen gleich hoch sind (Fliesstext), faellt das nicht auf. Sobald die
% Hoehen wechseln — Pillen einer Zielkette, Displaymathe, ein Bild —, wird der
% Rhythmus zum Nebenprodukt der Tinte: gemessen an der Zielkette der Showcase
% lagen zwischen benannt gleichen Elementen 4,1 bis 11,3 pt, obwohl der Autor
% genau zwei Werte gesetzt hatte.
%
% \nointerlineskip nimmt die Zeilenschaltung für den naechsten Block heraus.
% Danach ist der Abstand von der UNTERKANTE des vorigen Blocks zur OBERKANTE
% des naechsten genau #1 — unabhaengig davon, was in den beiden steht. Weil
% Hoehe und Tiefe einer Zeile der Tinte folgen (Oberlaenge, Unterlaenge), ist
% das zugleich der optische Abstand.
%
% Das `-\parskip` bleibt aus demselben Grund wie zuvor: Der folgende Block ist
% ein Absatz, und TeX setzt \parskip beim Absatzbeginn selbst. Ohne den Abzug
% käme er zu #1 hinzu.
%
% Die erste Zeile faengt den Fall ab, dass zwei Einschuebe direkt aufeinander
% treffen (\ZSFsep, unmittelbar gefolgt vom naechsten Kettenglied): Dann gilt
% der zweite, statt dass sich beide addieren.
\newcommand{\ZSFInterlude}[1]{%
  \par
  \ifdim\lastskip>0pt\relax \vskip-\lastskip \fi
  \vskip\dimexpr#1-\parskip\relax
  \nointerlineskip
}
% Trenner-Abstand um den Mittelpunkt im Index-Locator "6.2 · (S. 4)" (66_index).
% Bewusst als Makro statt \newlength: das em soll an der Verwendungsstelle in der
% Index-Schrift ausgewertet werden. Eine Länge würde es in der Präambel-Schrift
% einfrieren und den Abstand dadurch messbar verschieben.
\newcommand{\ZSFindexLocatorSep}{\hspace{0.2em}}

\input{styles/40_colors_structure}
% ============================================================
% 50_typography_semantics.tex — Schriftmakros für Box-Typen & semantische Inhaltsmakros
% ============================================================

% Explizite bp-Grössen (wie Analysis) für präzise, font-unabhängige Kontrolle
\newcommand{\ZSFfontChapter}{\sffamily\bfseries\fontsize{8.1bp}{9.1bp}\selectfont}
\newcommand{\ZSFfontSection}{\sffamily\bfseries\fontsize{8.1bp}{9.1bp}\selectfont}
\newcommand{\ZSFfontSubsubsection}{\sffamily\bfseries\fontsize{7.2bp}{8.0bp}\selectfont}
\newcommand{\ZSFfontBoxTitle}{\sffamily\bfseries\fontsize{8.1bp}{9.1bp}\selectfont}
% chktex 6
\newcommand{\ZSFfontNote}{\sffamily\fontsize{6.7bp}{7.5bp}\itshape\selectfont}
% Rechtsbuendiges Meta-/Kategorie-Tag im Titelbalken (vgl. \ZSFTitleTag in
% styles/60_boxes.tex). Etwas kleiner und kursiv, damit es den Titel begleitet
% statt mit ihm zu konkurrieren.
\newcommand{\ZSFfontTitleTag}{\sffamily\fontsize{6.4bp}{7.0bp}\itshape\selectfont}
\newcommand{\ZSFfontTitleHeader}{\sffamily\fontsize{11bp}{12.5bp}\bfseries\selectfont}
\newcommand{\ZSFfontFooter}{\sffamily\bfseries\fontsize{9.5pt}{9.5pt}\selectfont}
\newcommand{\ZSFfontIdentityMarker}{\fontsize{4pt}{4pt}\selectfont}

% Schrift/Farbe für Code-Kommentare — nur für das optionale Code-Modul
% (styles/67_code_comments.tex) benötigt. ce_green stammt aus
% styles/40_colors_structure.tex. Schadlos wenn das Modul deaktiviert ist.
\newcommand{\ZSFfontCodeComment}{\ttfamily\color{ce_green}}

% Semantische Inhalts-Fontgrössen — die zwei Stufen des Reglers `font`
% (`normal` | `dense`). Kapitel wählen sie nicht direkt, sondern über den
% Regler an der Box (styles/60_boxes.tex) oder an der ZSFtable
% (styles/20_tables.tex). `dense` ist die eine Stufe kleiner: lange Register,
% Vergleichsmatrizen, gedrängte Aufzählungen.
%
% EIN Paar für beide Bausteine, nicht je eines: »eine Stufe kleiner« ist
% dieselbe Entscheidung, ob sie eine Tabelle oder einen Boxinhalt trifft. Mit
% zwei Paaren müsste jede Änderung am einen im anderen nachgezogen werden
% (styles/README → Vereinigungstest).
%
% Beide holen die Display-Skips zurück: \normalsize und \footnotesize setzen
% sie auf die Klassenwerte, und die ZSF-Entscheidung dazu ist eine andere
% (\ZSFApplyDisplaySkips in 30_layout_spacing). Ohne diesen Nachsatz wären
% Formeln in jeder Box luftiger als im Fliesstext — lautlos, weil beides
% für sich korrekt aussieht.
\newcommand{\ZSFfontContent}{\normalsize\ZSFApplyDisplaySkips}
\newcommand{\ZSFfontContentDense}{\footnotesize\ZSFApplyDisplaySkips}

%
%
%
%
%
%
%
%
%
\newcommand{\ZSFfontTableHead}{\bfseries}
% Marke einer ZSFlist und Beschriftung eines \ZSFsep-Blockwechsels — beide
% sind dieselbe Rolle: eine kurze fette Beschriftung innerhalb eines Bausteins.
\newcommand{\ZSFfontBlockLabel}{\bfseries}
% Text, der per \textVorBox/\textNachBox an eine Box gebunden wird.
\newcommand{\ZSFfontBoxBinding}{\small}
% Inline-Pille für Stolperfallen (\ZSFdanger).
\newcommand{\ZSFfontDangerTag}{\sffamily\bfseries}
%
%
%
\newcommand{\ZSFfontRefLocator}{\upshape\bfseries}

%
%
%
%
%
\newcommand{\ZSFBoxNote}[1]{{\ZSFfontNote\color{MutedTextColor}#1}}

%
%
%
%
%
%
%
%
%
%
%
%
%
%
%
%
%
%
%
%
%
%
%
%
%
%
\providecommand{\ZSFDiagramLabelScale}{1}
% Die Grösse wird gerechnet statt geschrieben. Über \the und damit in pt, weil
% \fontsize sein Argument sonst als Zahl ohne Einheit lesen müsste — der Faktor
% kommt so als fertiges Mass an und nicht als Term, den \fontsize zerlegt.
\makeatletter
\newlength{\ZSF@fontSize}
\newlength{\ZSF@fontSkip}
\newcommand{\ZSF@scaledFont}[3]{%
  \setlength{\ZSF@fontSize}{#1\dimexpr #2\relax}%
  \setlength{\ZSF@fontSkip}{#1\dimexpr #3\relax}%
  \fontsize{\the\ZSF@fontSize}{\the\ZSF@fontSkip}\selectfont}
\newcommand{\ZSFfontDiagramLabel}{%
  \sffamily\ZSF@scaledFont{\ZSFDiagramLabelScale}{5.8bp}{6.4bp}}
\newcommand{\ZSFfontDiagramLabelSmall}{%
  \sffamily\ZSF@scaledFont{\ZSFDiagramLabelScale}{4.6bp}{5.1bp}}
\makeatother

%
%
%
%
%
%
%
%
%
%
%
%
%
%
%
%
%
%
%
%
%
%
%
%
\NewDocumentCommand{\ZSFkeyword}{s o m}{\ZSFkeywordStyle{#3}}

%
%
%
%
%
\newcommand{\ZSFlabel}[1]{\ZSFlabelStyle{#1}}

% Die beiden Darstellungs-Hooks der Marker — getrennt, weil \ZSFkeyword und
% \ZSFlabel verschiedene Dinge bedeuten und eine ZSF den Unterschied ändern
% können soll, ohne ein einziges Kapitel anzufassen.
%
% BEIDE schlicht fett, ohne Farbe. \ZSFkeyword stand bis hierher fett in der
% KAPITELFARBE, mit dem Argument, der primäre Scan-Anker müsse sich von jeder
% anderen Fettung abheben. Das Argument ist am Satz gescheitert: In laufendem
% Text bringt eine eingefärbte Vokabel Unruhe in den Absatz, ohne dass die
% Farbe etwas sagt — sie wiederholt nur, in welchem Kapitel man ohnehin gerade
% ist. Wer den Absatz liest, wird gestört; wer ihn überfliegt, findet die
% Fettung auch ohne Farbe.
%
% Farbe bleibt damit im System das, was sie vorher schon war und wofür sie
% trägt: Kapitel-Identität auf Flächen (Balken, Akzente), Wegweiser zum ZIEL
% eines Verweises (\ZSFref) und semantische Zuordnung in Formeln (\ZSFmhl*).
% Als blosse Betonung im Fliesstext ist sie ausdrücklich nicht vorgesehen.
%
% Die Unterscheidung der beiden Marker liegt jetzt allein in der Bedeutung —
% \ZSFkeyword ist ein Fachbegriff und landet im Register, \ZSFlabel ist eine
% neutrale Beschriftung und tut das nicht. Das ist der Unterschied, der beim
% Schreiben zählt; er war nie der Farbe zu entnehmen.
\newcommand{\ZSFkeywordStyle}[1]{\textbf{#1}}
\newcommand{\ZSFlabelStyle}[1]{\textbf{#1}}

\newcommand{\ZSFScriptRef}[1]{\hfill\textcolor{\ScriptRefColor}{\ZSFfontBoxBinding (S.#1)}}

% \ZSFhl{...} — Hervorhebung des EINEN prüfungskritischsten Fakts pro Box
% (fett + unterstrichen). Bewusst sparsam: hoechstens einmal pro Block.
% Klare Rollentrennung:
%   * Fachbegriff/Scan-Anker        -> \ZSFkeyword
%   * Stolperfalle/kritische Ausnahme -> \ZSFdanger
%   * Folgerung                      -> \ZSFconclusion
%   * sonst die EINE Schlüsselaussage -> \ZSFhl
% \uline statt \underline: \underline steckt seinen Inhalt in eine
% unzerbrechliche Box — jede Hervorhebung, die länger als eine Zeile ist,
% sprengt damit die 50-mm-Spalte. \ZSFhl markiert aber laut Zweck eine ganze
% Aussage, nicht ein Wort. \uline bricht mit dem Absatz um.
% Fett als Deklaration, nicht als Argument von \uline: ein \textbf{...} als
% Argument macht den Inhalt für ulem wieder zu einer Einheit.
%
% Die Darstellung liegt wie bei \ZSFkeyword und \ZSFlabel in einem eigenen
% Hook. Ohne ihn erwischte »die Hervorhebung dieser ZSF umstellen« zwei von
% drei Markern — der dritte war der einzige mit fest verdrahteter Erscheinung,
% und man sah es erst am Satz (styles/README → Auszeichnungs-Hooks).
\newcommand{\ZSFhlStyle}[1]{{\bfseries\uline{#1}}}
\newcommand{\ZSFhl}[1]{\ZSFhlStyle{#1}}

% Papierorientierte Referenzrollen:
%   * \ZSFref{label}: hilfreicher Querverweis mit farbiger Abschnittsnummer.
%   * \ZSFsectionref{label}: kompakte Zielnummer für lokale Router/Tabellen.
% Ziel-Label via optionales Argument von \SubsubsectionBar/\SubsectionBar/\StartChapter setzen:
%   \SubsectionBar[sec:matrix-inverse]{Inverse}
% Farbe = Palette-Farbe des ZIELkapitels (Wegweiser-Prinzip): wird automatisch
% aus der Kapitelnummer der Referenz abgeleitet ("6.6" -> ChapterColor6).
% Bei Kapiteln jenseits der Palette wird derselbe Modulo-Lookup wie bei
% \StartChapter genutzt, damit gewrappte Kapitel und Verweise konsistent sind.
% Aufrecht + fett statt kursiv für maximale Lesbarkeit bei 6.7bp.
\makeatletter
\def\ZSF@refExtractFirst#1#2\@nil{#1}
\def\ZSF@refExtractChap#1.#2\@nil{#1}
% Bestimmt \ZSF@refTargetColor aus dem Label; Fallback auf die aktive
% Kapitelfarbe im ersten latexmk-Pass (Label noch undefiniert).
\newcommand{\ZSF@refComputeTargetColor}[1]{%
  \@ifundefined{r@#1}{%
    \edef\ZSF@refTargetColor{\chaptercolor}%
  }{%
    \edef\ZSF@refTargetNum{\expandafter\expandafter\expandafter
      \ZSF@refExtractFirst\csname r@#1\endcsname\@nil}%
    \edef\ZSF@refTargetChap{\expandafter\ZSF@refExtractChap\ZSF@refTargetNum.\@nil}%
    \ZSFSetPaletteColorNameFromNumber{\ZSF@refTargetChap}{\ZSF@refTargetColor}%
  }%
}
% Die Zielfarbe läuft über \ZSFInk. Ohne das steht ein Verweis auf das eigene
% Kapitel im Titel einer \SubsectionBar als ChapterColorN auf \chaptercolor!75 —
% derselbe Farbton bei voller Sättigung, also unlesbar, ohne dass Build oder
% Linter etwas melden. Auf einer Fläche mit eigener Kontrastfarbe erbt der
% Verweis diese und bleibt lesbar; der Farb-Wegweiser gilt dort, wo er auch
% ankommt — im Boxinhalt (40_colors_structure → Ink-Vertrag).
\newcommand{\ZSFref}[1]{%
  \begingroup
  \ZSF@refComputeTargetColor{#1}%
  \hyperref[#1]{{\ZSFfontNote\ZSFfontRefLocator
    \ZSFInk{\ZSF@refTargetColor}{($\rightarrow$\,\ref*{#1})}}}%
  \endgroup
}
\newcommand{\ZSFsectionref}[1]{%
  \begingroup
  \ZSF@refComputeTargetColor{#1}%
  \hyperref[#1]{{\ZSFfontRefLocator\ZSFInk{\ZSF@refTargetColor}{\ref*{#1}}}}%
  \endgroup
}
% \ZSFhubref{label} — Verweis mit Zielnummer UND Zielname. [FORK-ERWEITERUNG]
%
% Dritte Rolle derselben Familie, nicht eine dritte Schreibweise derselben
% Sache. Der Unterschied ist, wer den Kontext trägt:
%   \ZSFref        Nummer neben einer Aussage, die selbst schon sagt, worum es
%                  geht — der Verweis ist Zusatz.
%   \ZSFsectionref blosse Zielnummer in einer Tabellenzelle, deren Zeile den
%                  Kontext trägt.
%   \ZSFhubref     der Verweis IST der Eintrag. In einer Übersicht, die auf
%                  mehrere Stellen zeigt, steht daneben nichts, was die Nummer
%                  erklären würde; \ZSFref ergäbe dort eine Reihe blinder
%                  Sprungziele.
%
% Den Klartext liefert \nameref von selbst: Die Strukturmakros hinterlegen den
% Titel bereits bereinigt (\ZSFplainName, oben). Dieses Makro braucht deshalb
% keine eigene Liste stillzulegender Marker — die stünde sonst hier ein zweites
% Mal und veraltete beim nächsten neuen Marker.
\newcommand{\ZSFhubref}[1]{%
  \begingroup
  \ZSF@refComputeTargetColor{#1}%
  \hyperref[#1]{{\ZSFfontNote\ZSFfontRefLocator
    \ZSFInk{\ZSF@refTargetColor}{($\rightarrow$\,\ref*{#1}\,\nameref*{#1})}}}%
  \endgroup\allowbreak
}

% ------------------------------------------------------------
% Der Klartext-Name eines Ziels — die Grundlage von \nameref
% ------------------------------------------------------------
% Ein Titel darf Marker tragen: \SubsectionBar[sec:inverse]{Inverse
% \ZSFref{sec:hub-full-rank}}. Fuer die ANZEIGE des Balkens ist das richtig;
% als NAME des Labels ist es falsch — ein \ZSFref darin ergaebe einen Verweis
% im Verweis, ein \ZSFkeyword einen Registereintrag an der Fundstelle des
% Lesers statt an der des Ziels.
%
% Die Bereinigung steht deshalb an EINER Stelle und wird beim Setzen des
% Labels eingebacken (\ZSF@storeLabelName, aufgerufen aus den Strukturmakros).
% Jeder Konsument von \nameref bekommt sie automatisch — auch ein handge-
% schriebenes \nameref im Kapitel. Die Alternative waere, dass jede Lesestelle
% dieselbe Liste mitfuehrt; beim naechsten neuen Marker fehlt sie dort.
%
% \ZSFplainName traegt bewusst KEIN @: Der Name wandert ueber die .aux-Datei,
% und die wird ohne \makeatletter gelesen — gleiche Regel wie bei den
% Index-Makros \ZSFidxnum/\ZSFidxsee.
\NewDocumentCommand{\ZSF@plainKeyword}{s o m}{#3}
% Die Schlucker nehmen ihr Argument UND das Leerzeichen davor: Steht der Marker
% am Titelende (»Inverse \ZSFref{…}«), bliebe sonst ein Leerzeichen stehen.
\newcommand{\ZSF@eatOne}[1]{\unskip}
\newcommand{\ZSF@eatTwo}[2]{\unskip}
\NewDocumentCommand{\ZSFplainName}{+m}{%
  \begingroup
    \let\ZSFkeyword\ZSF@plainKeyword
    \let\ZSFlabel\@firstofone
    \let\ZSFhl\@firstofone
    \let\ZSFref\ZSF@eatOne
    \let\ZSFsectionref\ZSF@eatOne
    \let\ZSFScriptRef\ZSF@eatOne
    \let\ZSFTitleTag\ZSF@eatOne
    \let\ZSFindex\ZSF@eatOne
    \let\ZSFindexsee\ZSF@eatTwo
    #1%
  \endgroup
}
% Hinterlegt den Titel als Namen des folgenden \label. Ohne das liefert
% \nameref dokumentweit nichts — lautlos, weil kein Baustein es bisher las.
\newcommand{\ZSF@storeLabelName}[1]{\def\@currentlabelname{\ZSFplainName{#1}}}
\makeatother

\newcommand{\ZSFconclusion}[1]{%
  \ZSFInterlude{\ZSFspaceS}%
  \noindent$\Rightarrow$ #1%
  \ZSFInterlude{\ZSFspaceS}%
}



% ── Marker-Darstellung ───────────────────────────────────────────────
% Default: \ZSFkeyword und \ZSFlabel beide schlicht fett, ohne Farbe.
% Inline-Betonung im Fliesstext trägt bewusst KEINE Kapitelfarbe — sie brächte
% Unruhe in den Absatz und sagte nichts, was die Seite nicht ohnehin zeigt.
% Farbe bleibt den Flächen, den Verweisen und den Formel-Markern vorbehalten.
% Umstellen heisst GENAU diese eine Zeile, nie ein Kapitel:
% \renewcommand{\ZSFkeywordStyle}[1]{\textbf{#1}}
% \renewcommand{\ZSFlabelStyle}[1]{\textbf{#1}}
% Dritter Marker derselben Familie: die Hervorhebung der Kernaussage.
% \renewcommand{\ZSFhlStyle}[1]{{\bfseries #1}}

% ── Grössenfarben: Farbe als Identität einer Grösse ──────────────────
% Für Fächer, in denen dieselben Grössen dokumentweit in Formeln auftauchen.
% LinAlg arbeitet stattdessen positional mit \ZSFmhlA..D innerhalb einer
% Herleitung — hier wird bewusst nichts deklariert.
% \ZSFDeclareQuantity{Kraft}{1}

% ── Eigener Ton: eine wiederkehrende Farbwelt fürs Fach ──────────────
% `tone` hat drei mitgelieferte Werte (chapter, neutral, warn).
% \ZSFDeclareTone{beispiel}{ ... }

% ============================================================
% 55_readability.tex — Intelligente Zeilenumbruch-Steuerung
% ============================================================
% Ziel: In schmalen Spalten (hier 4-spaltiges Landscape-A4) bricht
% der Satz lieber an natürlichen Wortgrenzen, statt die Zeile mit
% Leerraum aufzublähen oder Wörter hart zu trennen, nur um den
% rechten Rand zu erreichen.
%
% Kernbausteine:
%   \ZSFReadableParams : feingetunte TeX-Penalties (Flattersatz-kompatibel)
%   \ZSFReadableOn     : aktiviert RaggedRight + Parameter im aktuellen Scope
%   \ZSFReadableAuto   : respektiert den globalen Schalter
%   ZSFReadable        : Environment für lokale Anwendung
%   \ZSFbreak / \ZSFnobreak / \ZSFallowbreak : semantische Feinsteuerung
%   \ifZSFReadableBody : globaler Schalter (Standard: an)
%
% Dokumentweite Regel: Direkt nach \begin{multicols*} aktiviert
% \ZSFReadableAuto den Flattersatz für den gesamten Spalteninhalt.
% Lokale Anwendungen in Boxen bleiben erlaubt und sind unschädlich.
% ============================================================

% ------------------------------------------------------------
% Globaler Schalter: erlaubt das System projektweit aus-/einzuschalten
% ------------------------------------------------------------
\newif\ifZSFReadableBody
\ZSFReadableBodytrue
\newcommand{\ZSFReadableBodyOn}{\ZSFReadableBodytrue}
\newcommand{\ZSFReadableBodyOff}{\ZSFReadableBodyfalse}

% ------------------------------------------------------------
% TeX-Parameter für "lieber früh brechen als quetschen / hyphen"
% Werte bewusst moderat, damit Silbentrennung möglich bleibt, aber
% nicht die erste Wahl ist (wichtig in 4-Spalten-Layout).
% ------------------------------------------------------------
\newcommand{\ZSFReadableParams}{%
  \pretolerance=1200          % erst ohne Hyphenation versuchen (Default 100)
  \tolerance=2400             % dann Hyphenation erlauben (Default 200)
  \hyphenpenalty=200          % Silbentrennung leicht teurer -> Wortgrenze bevorzugt
  \exhyphenpenalty=200        %
  \doublehyphendemerits=10000 % zwei Bindestriche in Folge vermeiden
  \finalhyphendemerits=10000  % letzter Absatz-Zeile nicht hyphenisieren
  \emergencystretch=\ZSFreadabilityEmergencyStretch % Notfall-Dehnung gegen overfull hbox
  \hbadness=2000              % leise bei kleinen underfulls
  \linepenalty=10             % extra Zeilen sind günstig -> weniger Quetschen
}

% ------------------------------------------------------------
% \ZSFReadableOn: Flattersatz (RaggedRight aus ragged2e) + Parameter
%   - RaggedRight behält Silbentrennung als Fallback bei.
%   - rightskip hat etwas Flexibilität, damit kein "Waterfall"-Look entsteht.
% ------------------------------------------------------------
\newcommand{\ZSFReadableOn}{%
  % Schmale Flatterzone: natürliche Wortabstände, aber gut gefüllte Spalten.
  \setlength{\RaggedRightRightskip}{\ZSFreadabilityRaggedRightskip}%
  \RaggedRight
  \ZSFReadableParams
}

% Environment-Form für lokalen Einsatz
\newenvironment{ZSFReadable}{\par\ZSFReadableOn\ignorespaces}{\par}

% ------------------------------------------------------------
% Komfort-Wrapper: bedingtes Aktivieren (gesteuert über \ifZSFReadableBody)
% ------------------------------------------------------------
\newcommand{\ZSFReadableAuto}{\ifZSFReadableBody\ZSFReadableOn\fi}

%
%
%
%
%
%
%
%
%
%
%
%
%
%
%
%
%
%
%
%
\newcommand{\ZSFbreak}{\unskip\penalty-50\hskip\fontdimen2\font
  plus\fontdimen3\font minus\fontdimen4\font\relax}

% \ZSFnobreak — verhindert Umbruch an dieser Stelle (gebundenes Leerzeichen).
%   Nutzbar für Einheiten: "5\ZSFnobreak kg", "Punktladung\ZSFnobreak q_1"
\newcommand{\ZSFnobreak}{\nobreak\ }

% \ZSFallowbreak — erlaubt Umbruch an einer sonst FESTEN Stelle, also dort, wo
%   kein Leerzeichen steht ("Wortgrenze\ZSFallowbreak Fortsetzung" bleibt ein
%   Wort). Deshalb bringt es bewusst KEINEN Zwischenraum mit — das ist der
%   Unterschied zu \ZSFbreak und nicht eine vergessene Zeile.
%
%   \penalty0 statt \penalty100: »erlaubt« heisst neutral. Eine positive Strafe
%   erlaubt den Umbruch zwar auch, macht ihn aber teurer als jede andere
%   Stelle — der Name sagte damit etwas anderes als der Wert tat.
\newcommand{\ZSFallowbreak}{\penalty0\relax}

%
%
%
%
%
%
%
%
%
%
%
%
%
%
%
%
%
%
%
%
%
%
%
%
%
\newcommand{\ZSFcolonBindPenalty}{9000}
\makeatletter
\newcommand{\ZSF@bind}[1]{%
  \ifmmode~\else\unskip\penalty#1\hskip\fontdimen2\font\relax\fi
}
% Die beiden oeffentlichen Marken MUESSEN in derselben \makeatletter-Klammer
% stehen: \ZSF@bind ist ein Kontrollwort mit @, und nach \makeatother zerfiele
% es beim Tokenisieren in \ZSF und @bind.
\newcommand{\ZSFcolonbind}{\ZSF@bind{\ZSFcolonBindPenalty}}
\newcommand{\ZSFunitbind}{\ZSF@bind{10000}}
\makeatother

% ------------------------------------------------------------
% Automatischer LuaLaTeX-Filter für intelligente Zeilenumbrüche
% ------------------------------------------------------------
% Lädt den Lua-Filter zur Vermeidung von Trennungen bei Doppelpunkten,
% Einheiten und kurzen Inline-Formeln. Die Registrierung erfolgt
% über \AtBeginDocument, damit Paketladevorgänge ungestört bleiben.
\directlua{dofile(kpse.find_file("styles/zsf_filter.lua"))}
\AtBeginDocument{%
  \directlua{zsf_register_filter()}%
}

% ============================================================


% ── Fliesstext: Flattersatz ──────────────────────────────────────────
% \ZSFReadableBodyOn ist der Default, weil die ~50 mm schmalen Spalten im
% Blocksatz entweder trennen oder Wortzwischenräume aufblähen müssen.
% \ZSFReadableBodyOff             % zurück zu Blocksatz

% ============================================================
% 60_boxes.tex — tcolorbox-Stile, Balken, Boxen, Formelnoten, Listen
% ============================================================

% ------------------------------------------------------------
% Balken-Bindung: ein Titelbalken steht nie allein am Spaltenfuss
% ------------------------------------------------------------
% Ein Balken ohne den Anfang seines Inhalts ist der auffälligste Layoutfehler
% einer ZSF: Man sieht die Überschrift, aber nicht, wozu sie gehört.
%
% Eine Platzreserve allein verhindert das nicht, sie verursacht es: Sie besteht
% aus \penalty0 … \penalty-100 und legt damit einen BILLIGEN Umbruchpunkt an
% genau die Stelle, an der ein Balken gerade fertig gesetzt wurde. Ihre Bremse
% — Stretch-Glue, die einen Umbruch bei viel Restplatz unattraktiv macht —
% trägt unter \raggedcolumns nicht, weil multicols den Spaltenboden ohnehin
% frei auffüllt. Übrig bliebe die blosse Einladung zum Umbruch.
%
% Deshalb zwei getrennte Aufgaben statt einer:
%   \ZSFReserveSpace  — „lieber vorher umbrechen, wenn es eng wird" (Bremse)
%   \ZSFTitleBarLock  — „hier NICHT umbrechen, es kam noch kein Inhalt" (Sperre)
% Die Sperre ist die verlässliche; sie hängt an keiner Höhenschätzung.
\newif\ifzsfTitleBarPending
\newcommand{\ZSFTitleBarLockOn}{\global\zsfTitleBarPendingtrue}
\newcommand{\ZSFTitleBarLockOff}{\global\zsfTitleBarPendingfalse}

% Jeder Umbruchpunkt, den das Balken-System selbst setzt, läuft hierüber:
% solange noch kein Inhalt folgte, wird daraus ein Verbot.
\newcommand{\ZSFBarPenalty}[1]{%
  \ifzsfTitleBarPending\penalty10000\else\penalty#1\fi
}

\newcommand{\ZSFReserveSpace}[1]{%
  \par
  \ifzsfTitleBarPending
    \penalty10000
  \else
    \vspace{#1}\penalty0\vspace{-#1}%
  \fi
  \relax
}
\newcommand{\BreakIfNotEnoughSpace}{\ZSFReserveSpace{\ZSFbreakThreshold}}
\newcommand{\BreakIfNotEnoughSpaceChapter}{\ZSFReserveSpace{\ZSFbreakThresholdChapter}}
\newcommand{\BreakIfNotEnoughSpaceSubsubsection}{\ZSFReserveSpace{\ZSFbreakThresholdSubsubsection}}

% Einheitlicher Top-Spacing-Kontrakt für Titelbalken.
% Reihenfolge ist kritisch: \addvspace MUSS vor \BreakIfNotEnoughSpace stehen,
% damit es per Max-Semantik mit dem after-skip der Vorbox kollabiert.
% tcolorbox-Definitionen, die diesen Helper rufen, MUESSEN before skip=0pt
% setzen, damit das Spacing nicht doppelt eingerechnet wird.
\newcommand{\ZSFTitleBarTopSpacing}[1]{%
  \par
  \addvspace{#1}%
  \BreakIfNotEnoughSpace
}

%
%
%
%
%
%
%
%
%
%
%
%
%
%
%
%
%
%
%
%
%
%
%
%
%
%
\newcommand{\ZSFTitleBarAfter}{%
  \par\penalty10000\ZSFInterlude{\ZSFbarAfterGap}\ZSFTitleBarLockOn
}

\tcbset{
  zsfbarcolor/.style={colback=#1, colframe=#1},
  zsfbar/.style={
    enhanced jigsaw,
    breakable=false,
    boxrule=0pt,
    boxsep=0pt,
    arc=\ZSFboxArc,
    outer arc=\ZSFboxArc,
    left=\ZSFboxPadX,
    right=\ZSFboxPadX,
    before skip=0pt,
  },
}

\newtcolorbox{chapterbar}[1][]{
  zsfbar,
  zsfbarcolor=\ChapterBarColor,
  top=\ZSFbarPadY,
  bottom=\ZSFbarPadY,
  before skip=\ZSFspaceL,
  after={\ZSFTitleBarAfter},
  fontupper=\ZSFfontChapter,
  #1
}
% \ZSFInkOwned steht in ALLEN vier Balken direkt neben der Kontrastfarbe:
% Ein Balken ist eine gesaettigte Flaeche mit eigener Textfarbe und besitzt
% damit die Farbe (40_colors_structure -> Ink-Vertrag). Bewusst gleich
% behandelt, auch wo die Flaeche hell ist (Unterabschnitt) — das Kriterium
% ist "setzt eigene Kontrastfarbe" und nicht eine Einzelfallabwaegung, sonst
% muesste man es bei jeder Farbaenderung neu entscheiden.
\newcommand{\ChapterBar}[1]{
  \BreakIfNotEnoughSpaceChapter
  \begin{chapterbar}
    \textcolor{\ChapterBarTextColor}{\ZSFInkOwned #1}
  \end{chapterbar}
}

% ------------------------------------------------------------
% Dokument-Titel-Masthead — einmalig vor dem Front-Chapter
% ------------------------------------------------------------
\newtcolorbox{zsftitleheader}[1][]{
  zsfbar,
  zsfbarcolor=ZSFDocumentTitleBarColor,
  left=\ZSFboxPadXBar,
  right=\ZSFboxPadXBar,
  top=\ZSFbarPadY,
  bottom=\ZSFbarPadY,
  after={\ZSFTitleBarAfter},
  halign=center,
  #1
}
\newcommand{\ZSFTitleHeader}[2]{%
  \ZSFBarPenalty{-500}%
  \begin{zsftitleheader}
    {\color{ZSFDocumentTitleTextColor}\ZSFInkOwned\ZSFfontTitleHeader #1\par}%
    \vspace{\ZSFspaceXS}%
    {\color{ZSFDocumentTitleBylineColor}\ZSFInkOwned\ZSFfontNote von~#2\par}%
  \end{zsftitleheader}
}

% ------------------------------------------------------------
% Subsection-Bar (+ unnummerierte Star-Variante)
% ------------------------------------------------------------
\newtcolorbox{subsectionbar}[1][]{
  zsfbar,
  zsfbarcolor=\SubsectionBarColor,
  top=\ZSFbarPadY,
  bottom=\ZSFbarPadY,
  after={\ZSFTitleBarAfter},
  fontupper=\ZSFfontSection,
  #1
}
\newtcolorbox{subsubsectionbar}[1][]{
  zsfbar,
  zsfbarcolor=\SubsubsectionBarColor,
  top=\ZSFbarPadYTight,
  bottom=\ZSFbarPadYTight,
  after={\ZSFTitleBarAfter},
  fontupper=\ZSFfontSubsubsection,
  #1
}
\makeatletter
% Pro Ebene EIN Rumpf; die Varianten liefern nur das Nummernpräfix (leer bei
% der Sternvariante).
%   #1 = Label (leer erlaubt)   #2 = Nummernpräfix   #3 = Titel
\newcommand{\ZSF@subsectionBarBody}[3]{%
  \ZSFTitleBarTopSpacing{\ZSFsubsectionBarBeforeSkip}%
  \ZSFBarPenalty{-500}%
  \ZSFCollapseAfterChapterBar
  \ZSF@storeLabelName{#3}%
  \if\relax\detokenize{#1}\relax\else\label{#1}\fi
  \begin{subsectionbar}
    \textcolor{\SubsectionBarTextColor}{\ZSFInkOwned #2#3}
  \end{subsectionbar}
}
\newcommand{\ZSF@subsubsectionBarBody}[3]{%
  \BreakIfNotEnoughSpaceSubsubsection
  \par\addvspace{\ZSFsubsubsectionBarBeforeSkip}%
  \ZSFBarPenalty{-300}%
  \ZSF@storeLabelName{#3}%
  \if\relax\detokenize{#1}\relax\else\label{#1}\fi
  \begin{subsubsectionbar}
    \textcolor{\SubsubsectionTitleColor}{\ZSFInkOwned #2#3}
  \end{subsubsectionbar}
}

\newcommand{\SubsectionBar}{\@ifstar{\SubsectionBarStar}{\SubsectionBarNumbered}}
\newcommand{\SubsectionBarNumbered}[2][]{%
  \refstepcounter{subsection}\setcounter{subsubsection}{0}%
  \ZSF@subsectionBarBody{#1}{\thesubsection. }{#2}}
\newcommand{\SubsectionBarStar}[2][]{\ZSF@subsectionBarBody{#1}{}{#2}}

\newcommand{\SubsubsectionBar}{\@ifstar{\SubsubsectionBarStar}{\SubsubsectionBarNumbered}}
\newcommand{\SubsubsectionBarNumbered}[2][]{%
  \refstepcounter{subsubsection}%
  \ZSF@subsubsectionBarBody{#1}{\thesubsubsection. }{#2}}
\newcommand{\SubsubsectionBarStar}[2][]{\ZSF@subsubsectionBarBody{#1}{}{#2}}

\makeatother

%
%
%
%
\newcommand{\ZSFNewColumn}{\columnbreak}

%
%
%
%
%
%
%
%
%
%
%
%
%
%
%
%
%
%
%
%
%
%
%
%
%
%
%
%
%
%
%
\tcbset{zsftitlebar/.style={
  fonttitle=\ZSFfontBoxTitle\ZSFInkOwned\strut,
  halign title=center,
  title style={minimum height=\ZSFtitleBarHeight},
  toptitle=\ZSFspaceXS,
  bottomtitle=\ZSFspaceXS,
  before title={},
  after title={},
}}

% Rechtsbuendiges Meta-/Kategorie-Tag im Titelbalken einer Titelbox.
% Verwendung im Titel-Argument der Box, z.B.:
%   \begin{defbox}[Quicksort\ZSFTitleTag{O(n log n)}] ... \end{defbox}
%   \begin{codebox}[Algorithmus\ZSFTitleTag{Python}]  ... \end{codebox}
% Schiebt das Tag per \hfill an den rechten Titelrand. Schrift: \ZSFfontTitleTag.
\newcommand{\ZSFTitleTag}[1]{\hfill{\ZSFfontTitleTag #1}}

% Zum "before"-Schlüssel in allen Box-Stilen unten: das führende \par ist Pflicht.
% "before" ERSETZT die Default-Aktion von tcolorbox (die selbst ein \par enthält) —
% es hängt nicht an. Ohne \par wird eine Box, vor der roher Fliesstext ohne
% Leerzeile steht, inline an den offenen Absatz gehängt: sie bekommt keine
% Spaltenbreite und ragt sichtbar in die Nachbarspalte. Im vertikalen Modus ist
% \par ein No-op, der Zusatz ist also folgenlos, wo ohnehin alles korrekt war.
% Gemeinsamer Grundvertrag aller Inhaltsboxen. Was hier steht, gilt fuer jede
% Box — die einzelnen Stile setzen nur noch ihre Unterschiede. Besonders das
% before-Hook ist Pflicht (Begruendung im Kommentar oben) und darf deshalb
% nicht pro Box wiederholt werden, wo man es vergessen kann.
% Schicht 1 — Fluss. Auch rahmenlose Boxen (weight=quiet, splitbox) laufen
% hierüber, damit Abstände und Balken-Bindung überall identisch greifen.
% DER ABSTAND VOR EINEM BLOCK — eine Stelle, und sie muss es sein.
%
% Hier standen `before skip` und `before` nebeneinander. In tcolorbox ist
% `before skip` aber KEIN eigener Schluessel, sondern ein Stil, der `before`
% schreibt (tcolorbox.sty: `before skip/.style={before={...}}`). Beide
% schreiben also dasselbe Makro, und der spaetere gewinnt — `before` stand
% hinten, damit war \ZSFBoxBeforeSkip wirkungslos. Nachgemessen: auf 40pt
% gesetzt aendert sich am Satz nichts.
%
% Die Folgen reichten weiter als der eine Wert: Der Abstand vor JEDER Box kam
% damit allein aus \parskip und gehoerte keinem Register. Die Kollaps-Mechanik
% nach einem Balken (damals \ifzsfAfterSubsectionBar) rechnete einen Wert aus, den
% niemand las, und \ZSFboxgroup schloss seine Zwischenraeume nicht wirklich.
%
% Deshalb EIN `before`, das beides tut. Die Skip-Logik ist die von tcolorbox
% selbst — dieselben Faelle (nichts auf der Liste, vertikaler Modus), damit
% hier keine zweite Auslegung derselben Frage entsteht. Gelesen wird
% \ZSFBoxBeforeSkip VOR dem Verbrauch des Merkers: Der Merker waehlt den Wert.
%
% Den Minipage-Sonderfall von tcolorbox braucht es nicht: \addvspace prueft
% LaTeXs eigenes @minipage selbst und tut am Anfang einer Minipage nichts —
% genau das, was jener Zweig von Hand herstellt. Ihn nachzubauen hiesse,
% \iftcb@minipage zu lesen, und das existiert ausserhalb einer tcolorbox gar
% nicht (gemessen: undefined im runintext).
\makeatletter
\newcommand{\ZSF@blockBefore}{%
  \ifnum\lastnodetype=-1\relax\else
    \par
    \ifvmode\addvspace{\glueexpr\ZSFBoxBeforeSkip-\parskip\relax}\fi
  \fi
  % BEIDE Merker, nicht nur der eine. Ein Block zwischen Kapitelbalken und
  % Abschnittsbalken heisst, dass die beiden Balken NICHT aneinanderstossen —
  % blieb der Kapitel-Merker stehen, zog der Abschnittsbalken seinen Vorabstand
  % trotzdem zusammen. Derselbe Fehler eine Ebene hoeher.
  \ZSFConsumeAfterChapterBar\ZSFTitleBarLockOff
}
\tcbset{zsfbox/.style={
  enhanced jigsaw,
  after skip=\ZSFBoxAfterSkip,
  before={\ZSF@blockBefore},
}}
\makeatother

%
%
%
%
%
%
%
%
%
%
%

% Polsterung ist zweidimensional: waagrecht entscheidet über Kante und
% Balkenfreiraum, senkrecht über die Höhe. Getrennt, weil die Fälle sich
% tatsächlich kreuzen (Wertegitter: waagrecht bündig, senkrecht knapp).
%
% `padx` setzt den Abstand von Inhalt UND Titel zur Kante — es ist eine Kante
% und deshalb ein Wert (Begründung bei zsftitlebar oben). Für zwei der drei
% Werte steht deshalb NUR `left`/`right` da: Das sind Sammelschlüssel von
% tcolorbox (`left={lefttitle=#1, leftupper=#1, leftlower=#1}`), der Titel
% folgt also von selbst. Ihn danebenzuschreiben wäre derselbe Wert zweimal —
% und beim nächsten Eingriff einmal anders.
%
% `padx=none` ist der Fall, in dem Titel und Inhalt auseinandergehen MÜSSEN,
% und deshalb der einzige, der beide nennt: Die Box polstert nicht, weil ihr
% Inhalt es selbst tut — tablebox und valuegrid lassen ihre Farbfläche bis an
% den Rahmen laufen und halten den Rand in der äussersten Zelle
% (\ZSFtableEdgePad, siehe 20_tables). Der Titel hat keine Zelle und muss den
% Wert direkt lesen; ohne diese Zeile stünde er als einziges Element der Box
% an der blanken Kante.
\tcbset{
  padx/.is choice,
  padx/normal/.style={left=\ZSFboxPadX,    right=\ZSFboxPadX},
  padx/bar/.style   ={left=\ZSFboxPadXBar, right=\ZSFboxPadX},
  padx/none/.style  ={left=0pt,            right=0pt,
                      lefttitle=\ZSFtableEdgePad, righttitle=\ZSFtableEdgePad},
  pady/.is choice,
  pady/normal/.style={top=\ZSFboxPadY,    bottom=\ZSFboxPadY},
  pady/tight/.style ={top=\ZSFtextMinPad, bottom=\ZSFtextMinPad},
  pady/none/.style  ={top=0pt,            bottom=0pt},
}

%
%
%
%
%
%
%
%
%
%
%
%
%
%
%
%
\newenvironment{ZSFinset}{%
  \par\advance\leftskip\ZSFboxPadX\advance\rightskip\ZSFboxPadX
}{%
  \par
}

% Ausrichtung des Boxinhalts. Als Hook statt als direktem 'before upper',
% damit die Eintritts-Sequenz einer Box an EINER Stelle steht und die
% Ausrichtung sie nicht ersetzen muss.
\newcommand{\ZSFBoxAlign}{}
% Kontext-Flag der Box (heute nur die formulabox). Ebenfalls Teil derselben
% Sequenz — sonst bräuchte jede Box mit Kontext ein eigenes 'before upper'.
\newcommand{\ZSFBoxContext}{}
% Absatzabstand im Boxinhalt. Textboxen setzen ihn eigenständig (\ZSFspaceS),
% Container um Tabellen und Listen erben den Dokumentwert — dort gibt es keine
% Absätze, und ein eigener Wert würde nur die Liste verschieben. Als Regler
% statt als Auslassung, damit „erbt bewusst" und „vergessen" unterscheidbar sind.
\newcommand{\ZSFBoxParskip}{}
%
%
%
%
%
%
%
%
%
\newcommand{\ZSFBoxFont}{}
% Die Ausrichtung ERSETZT die Justierung, sie ergänzt sie nicht: `align=left`
% ist der Flattersatz für schmale Spalten (55_readability), `align=center` ist
% zentrierter Satz. Beides zusammen wäre kein zentrierter Flattersatz, sondern
% zentrierter Satz mit den Umbruch-Penalties des Flattersatzes — sichtbar an
% anders gebrochenen Formelzeilen.
\tcbset{
  align/.is choice,
  align/left/.style  ={/utils/exec={\renewcommand{\ZSFBoxAlign}{\ZSFReadableAuto}}},
  align/center/.style={/utils/exec={\renewcommand{\ZSFBoxAlign}{\centering}}},
  bodyparskip/.is choice,
  bodyparskip/box/.style={
    /utils/exec={\renewcommand{\ZSFBoxParskip}{\setlength{\parskip}{\ZSFspaceS}}}},
  bodyparskip/inherit/.style={
    /utils/exec={\renewcommand{\ZSFBoxParskip}{}}},
  font/.is choice,
  font/normal/.style={/utils/exec={\renewcommand{\ZSFBoxFont}{\ZSFfontContent}}},
  font/dense/.style ={/utils/exec={\renewcommand{\ZSFBoxFont}{\ZSFfontContentDense}}},
}
% Titel-Merker der leisen Fassung. Gefüllt wird er von der Fabrik, gesetzt vom
% Gewicht-Regler weiter unten; hier stehen nur die Deklarationen, damit sie vor
% beidem existieren.
%
% \ifZSF@boxQuiet trägt die Gewicht-Wahl. Sie steht schon fest, BEVOR die
% Optionsliste läuft — die Fabrik liest sie in einem Vor-Lauf aus dem Aufruf
% (siehe dort). Deshalb darf der Ton weiter unten sie bereits abfragen.
\makeatletter
\newcommand{\ZSF@boxTitleText}{}
\newif\ifZSF@boxHasTitle
\newif\ifZSF@boxQuiet
\makeatother
\newcommand{\ZSFBoxQuietTitle}{}
%
%
%
%
%
%
%
%
%
%
%
%
%
%
%
\newcommand{\ZSFBoxSetup}{\ZSFBoxFont\ZSFBoxParskip\ZSFBoxAlign\ZSFBoxContext}
\newcommand{\ZSFBoxEnter}{\ZSFBoxSetup\ZSFBoxQuietTitle}

%
%
%
%
%
%
%
%
%
%
%
%
%
%
%
%
%
%
\tcbset{zsfboxcontract/.style={
  zsfbox,
  tone=chapter,
  align=left,
  bodyparskip=box,
  font=normal,
  % Hooks zurücksetzen, damit eine verschachtelte Box die Einstellungen der
  % äusseren nicht erbt — sie sind gruppenlokal, aber nicht selbstlöschend.
  /utils/exec={\renewcommand{\ZSFBoxContext}{}\renewcommand{\ZSFBoxQuietTitle}{}},
  before upper={\ZSFBoxEnter},
  % Der Abstand zwischen oberer und unterer Haelfte, wenn sie gestapelt
  % stehen. Er gehoert in den Vertrag und nicht in ein Preset: JEDE Box kann
  % ein \tcblower tragen. Gesetzt hatte ihn nur die formulabox — alle
  % uebrigen bekamen tcolorbox' 2mm, ein Mass, das niemand gewaehlt hat und
  % das der Dichte nicht folgt. Derselbe Wert wie \ZSFsep: Es ist derselbe
  % Blockwechsel, nur von tcolorbox gezogen statt vom Autor.
  middle=\ZSFspaceSep,
  frame=soft,
  ZSF@surface=auto,
  ZSF@titlefill=tone,
  splitalign=top,
}}

%
%
%
%
%
%
%
%
%
%
%
%
%
%
%
%
%
%
%
%
%
%
\makeatletter
\newcommand{\ZSF@frameChoice}{soft}
\tcbset{
  frame/.is choice,
  frame/soft/.code  ={\renewcommand{\ZSF@frameChoice}{soft}},
  frame/strong/.code={\renewcommand{\ZSF@frameChoice}{strong}},
  frame/hard/.code  ={\renewcommand{\ZSF@frameChoice}{hard}},
  frame/none/.code  ={\renewcommand{\ZSF@frameChoice}{none}},
}
% Läuft als letzte Option jeder Box und übersetzt die gemerkte Stärke gegen
% den dann feststehenden Ton. Jede Box, die einen Rahmen haben kann, muss ihn
% durchlaufen — sonst bleibt sie beim tcolorbox-Default.
\tcbset{ZSF@resolveframe/.code={%
  \IfStrEqCase{\ZSF@frameChoice}{%
    {soft}  {\tcbset{colframe=\ZSFtoneFrameSoft}}%
    {strong}{\tcbset{colframe=\ZSFtoneFrameStrong}}%
    {hard}  {\tcbset{colframe=\ZSFtoneFrameHard}}%
    {none}  {\tcbset{frame hidden, boxrule=0pt, arc=0pt, outer arc=0pt}}%
  }%
}}
\makeatother

% Fläche — der Baustein wählt die ROLLE, die Farbe dazu kommt aus dem Ton.
%
% Wortgleich zum Rahmen darüber und aus demselben Grund: tcolorbox löst
% `colback` und `colbacktitle` beim Setzen der Option sofort per \colorlet auf.
% Beide gehören damit dem letzten Schreiber. Ein Preset, das seine Fläche
% direkt setzt, verliert sie deshalb an jede spätere Tonwahl des Aufrufers —
% und ein ausgeschriebenes `tone=chapter` sähe anders aus als ein
% weggelassenes, obwohl es dieselbe Wahl ist. Verzögert aufgelöst stehen
% Rolle und Ton beide fest, wenn der Resolver läuft; die Reihenfolge im
% Aufruf ist dann bedeutungslos.
%
%   auto      laut oder ruhig — entscheidet das Gewicht (Default)
%   plain     ungetönt; Basis für Zebra, Gitterlinien, Syntaxfarben
%   quiet     immer die ruhige Fläche, unabhängig vom Gewicht
%   emphasis  betont; für Boxen, deren Inhalt selbst die Aussage ist
%
% `ZSF@titlefill` beantwortet dieselbe Frage für die Kopfzeile: `tone` gibt
% ihr den Titelbalken des Tons, `surface` setzt sie auf die Fläche der Box
% (Formelbox — dort ist der Titel eine Zeile, kein Balken).
\makeatletter
\newcommand{\ZSF@surfaceRole}{auto}
\newcommand{\ZSF@titleFillRole}{tone}
\tcbset{
  ZSF@surface/.is choice,
  ZSF@surface/auto/.code    ={\renewcommand{\ZSF@surfaceRole}{auto}},
  ZSF@surface/plain/.code   ={\renewcommand{\ZSF@surfaceRole}{plain}},
  ZSF@surface/quiet/.code   ={\renewcommand{\ZSF@surfaceRole}{quiet}},
  ZSF@surface/emphasis/.code={\renewcommand{\ZSF@surfaceRole}{emphasis}},
  ZSF@titlefill/.is choice,
  ZSF@titlefill/tone/.code   ={\renewcommand{\ZSF@titleFillRole}{tone}},
  ZSF@titlefill/surface/.code={\renewcommand{\ZSF@titleFillRole}{surface}},
}
% Läuft als letzte Option jeder Box, direkt neben dem Rahmen-Resolver.
% Erst die Rolle in eine Farbe übersetzen, dann setzen — so kann die Kopfzeile
% dieselbe Farbe bekommen, ohne die Fallunterscheidung zu wiederholen.
\newcommand{\ZSF@resolvedBack}{\ZSFtoneLoudBack}
\tcbset{ZSF@resolvesurface/.code={%
  \IfStrEqCase{\ZSF@surfaceRole}{%
    {auto}    {\ifZSF@boxQuiet\renewcommand{\ZSF@resolvedBack}{\ZSFtoneQuietBack}%
               \else\renewcommand{\ZSF@resolvedBack}{\ZSFtoneLoudBack}\fi}%
    {plain}   {\renewcommand{\ZSF@resolvedBack}{\BoxPlainBackColor}}%
    {quiet}   {\renewcommand{\ZSF@resolvedBack}{\ZSFtoneQuietBack}}%
    {emphasis}{\renewcommand{\ZSF@resolvedBack}{\ZSFtoneEmphasisBack}}%
  }%
  \tcbset{colback=\ZSF@resolvedBack}%
  \IfStrEqCase{\ZSF@titleFillRole}{%
    {tone}   {}%
    {surface}{\tcbset{colbacktitle=\ZSF@resolvedBack}}%
  }%
}}
\makeatother

% Atomar: bricht nie über eine Spaltengrenze, auch nicht im Pack-Modus.
% Muss in der Optionsliste NACH dem Basisstil stehen — dorthin hängt
% \ZSFBoxesBreakableOn sein 'breakable'. Als benannter Regler ist diese
% Reihenfolge in den Presets festgeschrieben statt in einem Kommentar.
\tcbset{atomic/.style={unbreakable}}

% Schicht 2 — Form. Alles, was die Geometrie einer geschlossenen Inhaltsbox
% bestimmt. Wer Rahmenstärke, Ecken oder Innenabstand des Systems ändern will,
% ändert hier — und nur hier. padx/pady gehören dazu, damit \linewidth in jeder
% Box identisch ist (sonst variierende Tabellen-Breiten je Box-Typ).
%
% Was NICHT hier steht, sondern im Vertrag darüber: alles, was auch für eine
% Box mit eigener Form gilt (Ton, Justierung, Eintritt, Reset der verzögerten
% Wahlen). Sonst müsste jede rahmenlose Box es abschreiben.
\tcbset{zsfboxshape/.style={
  zsfboxcontract,
  breakable=false,
  boxsep=0pt,
  boxrule=\ZSFboxRuleWidth,
  titlerule=0pt,
  arc=\ZSFboxInnerArc,
  outer arc=\ZSFboxArc,
  padx=normal,
  pady=normal,
}}

%
%
%
%
%
%
%
%
%
%
%
%
%
%
%
%
%
%
%
%
%
%
%
%
%
%
%
%
%
%
%
%
%
%
%
%
%
%
%
%
%
%

% Schicht 3 — Titelbox. Der EINZIGE Vertrag für Boxen mit Titelbalken; alle
% Presets unten stehen darauf.
%
% Dass hier nur noch zwei Zeilen stehen, ist der Punkt: Form kommt aus
% zsfboxshape, alles Übrige aus dem Vertrag darüber. Eine Titelbox ist damit
% genau das, was ihr Name sagt — eine Box mit Titelbalken.
\tcbset{zsftitlebox/.style={
  zsfboxshape,
  zsftitlebar,
}}

%
%
%
%
%
%
%
%
%
\makeatletter
% Die Fallunterscheidung sitzt in einem tcolorbox-Schlüssel, nicht um das
% \begin herum: So gibt es genau EIN \begin{tcolorbox} und genau ein \end —
% sonst zählt chktex sie gegeneinander auf und meldet sie bis zum Dateiende
% als unbalanciert (Warnung 15, dort nicht mehr lokal unterdrückbar).
\tcbset{zsftitle/.code={\IfBlankTF{#1}{\tcbset{notitle}}{\tcbset{title={#1}}}}}
%
%
%
%
%
%
%
%
%
%
%
%
%
%
%
%
%
%
%
%
%
%
%
%
%
%
%
%
%
%
%
%
%
\newcommand{\ZSF@ownedKeyError}[2]{%
  \PackageError{ZSF}{%
    Schluessel '#1' gehoert dem Box-System, nicht dem Aufruf%
  }{%
    Flaeche und Rahmen werden am Ende jeder Optionsliste aus Rolle und Ton
    zusammengesetzt (ZSF@resolveframe / ZSF@resolvesurface). Ein '#1' im
    Aufruf wird dabei ueberschrieben und bliebe wirkungslos -- frueher
    stillschweigend, jetzt hier. Stattdessen den Regler '#2' verwenden;
    fehlt dafuer ein Wert, ist das eine Luecke und gehoert gemeldet
    (fehlender Optionswert).%
  }%
}
% Box-spezifische Regler gehören genau einer Box. Auf jeder anderen nimmt
% tcolorbox sie klaglos entgegen, und niemand liest sie je wieder — dieselbe
% stille Klasse wie die Farbschlüssel oben, nur eine Ecke weiter:
% `\begin{defbox}[T][ordered]` setzte die Listen-Nummerierung, ohne dass eine
% Liste existiert; `[grid=none]` leerte die Gitterwahl einer Box ohne Gitter.
%
% Welche Box welchen Regler besitzt, sagt sie selbst, bevor sie die Fabrik
% ruft; leer heisst »keinen«. Die Kommas sind Begrenzer, damit ein Reglername,
% der Teil eines anderen ist, nicht versehentlich durchgeht.
\newcommand{\ZSF@boxOwnKnobs}{}
\newcommand{\ZSF@requireOwnKnob}[2]{%
  \IfSubStr{\ZSF@boxOwnKnobs}{,#1,}{}{%
    \PackageError{ZSF}{%
      Regler '#1' gehoert zu #2, nicht zu dieser Box%
    }{%
      Auf dieser Box hat '#1' nichts, worauf er wirken koennte -- er wuerde
      stillschweigend verschwinden. Entweder die Box verwenden, der er
      gehoert (#2), oder die Option weglassen.%
    }%
  }%
}
\pgfkeys{
  /zsf/boxpre/.is family, /zsf/boxpre,
  weight/.is choice,
  weight/loud/.code ={\ZSF@boxQuietfalse},
  weight/quiet/.code={\ZSF@boxQuiettrue},
  ordered/.code={\ZSF@requireOwnKnob{ordered}{ZSFlist}},
  grid/.code   ={\ZSF@requireOwnKnob{grid}{valuegrid}},
  part/.code   ={\ZSF@requireOwnKnob{part}{codebox}},
%
%
%
%
%
%
  colback/.code     ={\ZSF@ownedKeyError{colback}{tone}},
  colframe/.code    ={\ZSF@ownedKeyError{colframe}{frame}},
  colbacktitle/.code={\ZSF@ownedKeyError{colbacktitle}{tone}},
  coltitle/.code    ={\ZSF@ownedKeyError{coltitle}{tone}},
  .unknown/.code={},% chktex 26
}
\NewDocumentEnvironment{ZSF@box}{m m m}{%
  \renewcommand{\ZSF@boxTitleText}{#2}%
  \IfBlankTF{#2}{\ZSF@boxHasTitlefalse}{\ZSF@boxHasTitletrue}%
  % Zentraler Reset statt einer Zeile pro Basisstil: Sonst erbt eine Box, deren
  % Stil den Merker nicht zurücksetzt, die Titelzeile der Box, in der sie steht.
  \renewcommand{\ZSFBoxQuietTitle}{}%
  \ZSF@boxQuietfalse
  \pgfkeys{/zsf/boxpre, #3}%
  % Der Besitz gilt für GENAU diese Box. Zurückgesetzt direkt nach dem
  % Vor-Lauf, damit eine verschachtelte Box ihn nicht erbt und dort ein
  % fremder Regler wieder stumm durchginge.
  \renewcommand{\ZSF@boxOwnKnobs}{}%
  \begin{tcolorbox}[#1, zsftitle={#2}, ZSF@weightBase, #3,
                    ZSF@resolveframe, ZSF@resolvesurface, ZSF@resolvesplit]%
}{%
  \end{tcolorbox}%
}
\makeatother

% ------------------------------------------------------------
% Presets — jede Box ist eine Liste aus den Reglern oben
% ------------------------------------------------------------
% Prüffrage an jedes Preset: Steht hier ein Schlüssel, für den es einen Regler
% gäbe? Dann gehört der Regler hin. Rohe Schlüssel bleiben nur, wo die
% Eigenschaft wirklich einmalig ist (der Zebra-Beschnitt der Tabelle, die harte
% Kontur der Formelbox) — und dann mit Begründung.

% Tabelle: Hintergrund bis an die Kante, damit das Zebra nicht kurz vor dem
% Rahmen endet. 'clip upper' schneidet die Streifen an der Rundung ab.
% Die Fläche bleibt ungetönt, weil das Zebra sie als Basis braucht.
\tcbset{preset/table/.style={
  zsftitlebox,
  ZSF@surface=plain,
  padx=none, pady=none,
  bodyparskip=inherit,
  clip upper,
}}
%
%
\tcbset{preset/fig/.style={zsftitlebox, atomic}}

% ------------------------------------------------------------
% Gewicht — laut oder leise, dieselbe Box
% ------------------------------------------------------------
% Laut oder leise ist eine Eigenschaft derselben Box, keine zweite Box:
%
%   weight=loud   Titelbalken, Rahmen, getönte Fläche  (Default)
%   weight=quiet  linker Akzentbalken, rahmenlos, ruhige Fläche
%
% Der Titel wandert dabei aus der Kopfzeile in die erste Inhaltszeile. Den
% Titeltext liefert die Fabrik (siehe dort), deshalb gilt der Regler auf JEDER
% Box und nicht nur auf der defbox.
\makeatletter
%
%
%
%
%
%
%
%
%
\newcommand{\ZSF@emitQuietTitle}{%
  \ifZSF@boxHasTitle
    {\ZSFfontBoxTitle\color{\ZSFtoneAccent}\ZSFInkOwned\ZSF@boxTitleText}%
    \par\vspace{\ZSFspaceXS}%
  \fi
}
% Die Basis der leisen Fassung. Sie wird von der Fabrik VOR den
% Aufrufer-Optionen eingesetzt (Begründung dort), nicht mittendrin — deshalb
% darf sie Vorbelegungen wie `frame=none` oder `padx=bar` mitbringen, ohne
% eine Wahl des Aufrufers zu verwerfen.
%
% Die Fläche setzt sie bewusst NICHT selbst: Sie kommt vom Ton, der beide
% Fassungen kennt und anhand des Gewichts wählt (siehe `ZSF@toneSurface`).
% Täte sie es hier, wäre sie wieder ein Schreiber auf `colback` neben dem Ton
% — und welcher gewinnt, hinge erneut an der Reihenfolge.
\tcbset{ZSF@weightBase/.code={%
  \ifZSF@boxQuiet
    \tcbset{preset/quiet, notitle,
            /utils/exec={\renewcommand{\ZSFBoxQuietTitle}{\ZSF@emitQuietTitle}}}%
  \fi
}}
% Der Schlüssel selbst bleibt als No-op bestehen: Er steht weiterhin in der
% Aufrufer-Liste, wo tcolorbox ihn sonst als unbekannt meldete. Gelesen wurde
% er zu diesem Zeitpunkt längst — im Vor-Lauf der Fabrik.
\tcbset{
  weight/.is choice,
  weight/loud/.code ={},
  weight/quiet/.code={},
}
\makeatother

\NewDocumentEnvironment{defbox}{O{} O{}}
  {\begin{ZSF@box}{zsftitlebox}{#1}{#2}}{\end{ZSF@box}} % chktex 9

\NewDocumentEnvironment{tablebox}{O{} O{}}
  {\begin{ZSF@box}{preset/table}{#1}{#2}}{\end{ZSF@box}} % chktex 9
% figbox und splitbox tragen einen Block — dort gilt fuer ein \ZSFimage das
% Budget einer Abbildung, nicht das der Tabellenzelle (siehe \ZSF@imgBudget).
% \makeatletter ist noetig, weil \ZSF@imgBudget ein Kontrollwort mit @ ist;
% \begin{ZSF@box} daneben braucht es nicht — ein Umgebungsname ist Text.
\makeatletter
\NewDocumentEnvironment{figbox}{O{} O{}}
  {\renewcommand{\ZSF@imgBudget}{\ZSFfigMaxHeight}%
   \begin{ZSF@box}{preset/fig}{#1}{#2}}{\end{ZSF@box}} % chktex 9
\makeatother
%
%
%
%
%
%
%
%
%
%
%
%
%
%
\NewDocumentEnvironment{warnbox}{O{Achtung} O{}}
  {\begin{ZSF@box}{zsftitlebox, tone=warn}{#1}{#2}}{\end{ZSF@box}} % chktex 9

%
%
%
%
%
%
%
%
%
%
%
%
%
%
%
%
%
%
%
%
%
%
%
%
%
%
%
%
%
%
%
%
%
%
%
%
\makeatletter
\newcommand{\ZSFgroupcols}{}
\newcommand{\ZSF@groupHead}{}
\newlength{\ZSF@groupOuterAfter}
\pgfkeys{
  /zsf/boxgroup/.is family, /zsf/boxgroup,
  cols/.store in=\ZSFgroupcols,
  head/.store in=\ZSF@groupHead,
  ZSF@groupDefaults/.style={cols=, head=},
}
\NewDocumentEnvironment{ZSFboxgroup}{O{}}{%
  \par
  \setlength{\ZSF@groupOuterAfter}{\ZSFBoxAfterSkip}%
  \addvspace{\ZSFBoxBeforeSkip}%
  \begingroup
  \pgfkeys{/zsf/boxgroup, ZSF@groupDefaults, #1}%
  % Innen keine Abstände: erst dadurch liest sich der Stapel als ein Block.
  \renewcommand{\ZSFBoxBeforeSkip}{0pt}%
  \setlength{\ZSFBoxAfterSkip}{0pt}%
  \ifx\ZSF@groupHead\@empty\else
    \begin{tablebox}%
      \begin{ZSFtable}[zebra=false]{\ZSFgroupcols}%
        \ZSFheaderRow\ZSF@groupHead \\%
      \end{ZSFtable}%
    \end{tablebox}%
  \fi
}{%
  \endgroup
  \par\addvspace{\ZSF@groupOuterAfter}%
}
\makeatother

% ------------------------------------------------------------
% Formula-Box
% ------------------------------------------------------------
% Der Unterschied zur Titelbox ist bewusst und klein: Der Titel sitzt auf dem
% Boxhintergrund statt auf einem eigenen Balken, der Inhalt ist zentriert,
% und zwischen upper und lower liegt eine Trennlinie. Alles andere — Rahmen,
% Ecken, Innenabstand, Titelschrift — kommt aus zsftitlebox.
%
% \ifZSF@inFormulaBox (deklariert in 30_layout_spacing) schaltet
% \textVorBox/\textNachBox auf ihre Formelform um. Das Flag wird in
% 'before upper' gesetzt und ist damit auf die Boxgruppe beschränkt — es
% braucht kein Zurücksetzen.
\makeatletter
\tcbset{preset/formula/.style={
  zsftitlebox,
  align=center,
  ZSF@surface=emphasis,
  ZSF@titlefill=surface,
  frame=hard,
  % Der waagrechte Titelabstand stand hier, weil er in zsftitlebar auf 0pt
  % festgenagelt war und auf der Fläche der Formelbox besonders auffiel. Er
  % kommt jetzt aus `padx` und gilt damit auf jeder Box — der Unterschied der
  % Formelbox ist allein, dass ihr Titel keine eigene Fläche hat.
  bottomtitle=0pt,
  /utils/exec={\renewcommand{\ZSFBoxContext}{\ZSF@inFormulaBoxtrue}},
  after upper={\par},
}}
\NewDocumentEnvironment{formulabox}{O{} O{}}
  {\begin{ZSF@box}{preset/formula}{#1}{#2}}{\end{ZSF@box}} % chktex 9
\makeatother

%
%
%
%
%
%
%
%
%
%
%
%
\newcommand{\ZSFBoxRule}{%
  \noindent\textcolor{FormulaSeparatorColor}{\rule{\linewidth}{\ZSFboxRuleWidth}}%
}

%
%
%
%
%
%
%
%
%
%
%
%
%
%
%
%
%
%
\NewDocumentCommand{\ZSFsep}{o}{%
  \ZSFInterlude{\ZSFspaceSep}%
  \ZSFBoxRule
  \IfValueT{#1}{\ZSFInterlude{\ZSFspaceS}{\ZSFfontBlockLabel #1\par}}%
  \ZSFInterlude{\ZSFspaceSep}%
  \penalty-50%
}

% Dezente Anmerkung als eigene Zeile — für die Erläuterung unter einer Formel
% oder Formelgruppe. Für eine Anmerkung, die auf der Formelzeile bleiben soll,
% \ZSFformulaline; für einen Satz, der eng an die Formel bindet, \textNachBox.
\newcommand{\formulanote}[1]{%
  \ZSFInterlude{\ZSFspaceS}%
  \noindent\ZSFBoxNote{#1}%
  \ZSFInterlude{\ZSFspaceS}%
}
% \ZSFformulaline{<mathe>}{<Anmerkung>} — Formelzeile mit rechtsbündiger
% Kurz-Anmerkung auf derselben Höhe (Gültigkeitsbereich, Einheit, Fallname).
\newcommand{\ZSFformulaline}[2]{%
  \par\noindent$\displaystyle #1$\hfill\ZSFBoxNote{#2}\par
}

% ------------------------------------------------------------
% Runintext + Grafik-Helpers
% ------------------------------------------------------------
% Der Fliesstext-Block nimmt seinen Vorabstand aus DERSELBEN Quelle wie jede
% Box (\ZSF@blockBefore). Vorher stand hier ein fester \ZSFspaceS: Nach einem
% Abschnittsbalken zog eine Box ihren Abstand zusammen, ein runintext nicht —
% derselbe Ort, zwei Abstände, je nachdem was folgt. Sichtbar war das als
% auffällige Lücke zwischen Balken und erstem Satz.
\makeatletter
\newenvironment{runintext}{%
  \ZSF@blockBefore
  \noindent\ZSFReadableAuto\ignorespaces
}{\par}
\makeatother

% Dezente Bildunterschrift (leerer Text -> nichts)
\newcommand{\ZSFfigcaption}[1]{%
  \if\relax\detokenize{#1}\relax\else
    \par\vspace{\ZSFspaceXS}%
    \ZSFBoxNote{#1}\par
  \fi
}

%
%
%
%
%
%
%
%
%
%
%
%
%
%
%
%
%
%
%
%
%
%
\makeatletter
\newcommand{\ZSF@figWidthFrac}{}
\newcommand{\ZSF@figBoxOpts}{}
\newlength{\ZSF@figHeightLimit}
% Oeffnet die figbox mit EXPANDIERTER Optionsliste. Ohne diesen Umweg landet
% der Makroname selbst in der Optionsliste: tcolorbox parst sie, bevor es
% expandiert, und meldet die ganze gesammelte Liste als einen unbekannten
% Schluessel. Der \expandafter am Aufrufort loest genau eine Stufe auf.
\newcommand{\ZSF@figOpen}[2]{\begin{figbox}[{#2}][#1]} % chktex 9
\pgfkeys{
  /zsf/fig/.is family, /zsf/fig,
  width/.store in=\ZSF@figWidthFrac,
  height/.code={\setlength{\ZSF@figHeightLimit}{#1}},
  .unknown/.code={% chktex 26
    \IfBlankTF{#1}%
      {\edef\ZSF@figBoxOpts{\ZSF@figBoxOpts,\pgfkeyscurrentname}}%
      {\edef\ZSF@figBoxOpts{\ZSF@figBoxOpts,\pgfkeyscurrentname=\unexpanded{#1}}}%
  },
  ZSF@figDefaults/.style={width=, height=\ZSFfigMaxHeight},
}

% ------------------------------------------------------------
% \ZSFimage — das Bild OHNE Box
% ------------------------------------------------------------
% \ZSFimage[width=<Anteil>, height=<Mass>]{pfad}
%
% Die Bild-Makros darüber liefern immer einen Block: figbox, Titel, Caption.
% Genau das ist falsch, sobald das Bild INHALT eines anderen Bausteins ist —
% in einer Tabellenzelle (Bauform neben Bezeichnung neben Kennwerten) gibt es
% keinen Platz für eine zweite Box. Bis hierher gab es dafür keinen Weg: rohes
% \includegraphics ist in Kapiteln verboten, \ZSFfig öffnet eine figbox.
%
% Zugleich die EINE Stelle im System, an der \includegraphics steht. Vorher
% dreimal fast gleich (Bildreihe, Einzelbild, Bild-neben-Text) — und wer die
% Skalierung ändern wollte, musste alle drei finden.
%
% Bewusst ohne \par und ohne \centering: Der Baustein, in dem das Bild steht,
% bestimmt Ausrichtung und Umbruch. Beides hier zu setzen machte das Makro in
% einer Tabellenzelle unbrauchbar — dem Fall, für den es existiert.
\newcommand{\ZSF@imgWidthFrac}{1}
\newlength{\ZSF@imgHeightLimit}
% Das Hoehenbudget gehoert dem CONTAINER, nicht dem Bild.
%
% Dasselbe \ZSFimage bedeutet an zwei Orten zwei Groessenordnungen: In einer
% Tabellenzelle muss es in eine Zeile passen (\ZSFimageMaxHeight), als
% Abbildung fuellt es einen Block (\ZSFfigMaxHeight). Das Bild selbst kann das
% nicht wissen — es sieht seinen Container nicht.
%
% Stuende hier ein fester Default, waere genau eine der beiden Lagen still
% falsch: Ein \ZSFimage in einer figbox wuerde ohne Fehler und ohne Warnung auf
% Zellenhoehe gestaucht, und die Hoehe muesste an jedem Aufruf wiederholt
% werden — also genau die Wiederholung, die ein zentraler Default vermeidet.
% Deshalb ist der Default eine INDIREKTION: Jeder Baustein, der Bilder tragen
% kann, bindet sie auf sein eigenes Mass um (figbox/splitbox auf Abbildungs-,
% ZSFtable/valuegrid auf Zellenhoehe). Wer nichts bindet, bekommt die Zelle —
% der Fall, fuer den \ZSFimage existiert.
\newcommand{\ZSF@imgBudget}{\ZSFimageMaxHeight}
\pgfkeys{
  /zsf/image/.is family, /zsf/image,
  width/.store in=\ZSF@imgWidthFrac,
  height/.code={\setlength{\ZSF@imgHeightLimit}{#1}},
  ZSF@imgDefaults/.style={width=1, height=\ZSF@imgBudget},
}
%
% Die senkrechte Luft in einer Tabellenzelle bringt das Bild NICHT selbst mit.
% Sie kommt von \ZSFtableCellPadY (20_tables): »Zellinhalt haelt Abstand zur
% Zeilenkante« gilt fuer ein Bild wie fuer einen Bruch, und der Regler `rows`
% entscheidet darueber. Ein eigener Bildwert waere dieselbe Aussage ein
% zweites Mal.
\newcommand{\ZSFimage}[2][]{%
  \begingroup
  \pgfkeys{/zsf/image, ZSF@imgDefaults, #1}%
  \includegraphics[width=\ZSF@imgWidthFrac\linewidth,
                   height=\ZSF@imgHeightLimit, keepaspectratio]{#2}%
  \endgroup
}

%
%
%
%
%
%
%
\newcounter{ZSF@figNum}
\newlength{\ZSF@figWidth}
\newcommand{\ZSFfig}[4][]{%
  \begingroup
  \renewcommand{\ZSF@figBoxOpts}{}%
  \pgfkeys{/zsf/fig, ZSF@figDefaults, #1}%
  \expandafter\ZSF@figOpen\expandafter{\ZSF@figBoxOpts}{#3}%
    \setcounter{ZSF@figNum}{0}%
    \@for\ZSF@figItem:=#2\do{\stepcounter{ZSF@figNum}}%
    \ifnum\value{ZSF@figNum}>1\relax
      \setlength{\ZSF@figWidth}{\dimexpr
        \ifx\ZSF@figWidthFrac\@empty\ZSFfigUsableFraction\else\ZSF@figWidthFrac\fi
        \linewidth/\value{ZSF@figNum}\relax}%
      \noindent
      \@for\ZSF@figItem:=#2\do{%
        \begin{minipage}[c]{\ZSF@figWidth}%
          \centering
          \ZSFimage[width=1, height=\ZSF@figHeightLimit]{\ZSF@figItem}%
        \end{minipage}%
        \hfill
      }%
      \unskip
      \par
    \else
      \centering
      % Vorbelegung erst hier aufloesen: In der Optionsliste von \ZSFimage
      % stuende sonst eine Fallunterscheidung, die pgfkeys als Wert speichert
      % statt sie auszuwerten.
      \ifx\ZSF@figWidthFrac\@empty\renewcommand{\ZSF@figWidthFrac}{1}\fi
      \ZSFimage[width=\ZSF@figWidthFrac, height=\ZSF@figHeightLimit]{#2}%
    \fi
    \ZSFfigcaption{#4}%
  \end{figbox}% chktex 10
  \endgroup
}

% \ZSFfigside[<Optionen>]{pfad}{Titel}{Inhalt rechts} — Bild links, Inhalt
% rechts. `width` ist hier der Anteil des BILDES (Vorbelegung 0.30); er wird
% erst am Einsatzort gegen die dortige \linewidth ausgewertet.
\newcommand{\ZSFfigside}[4][]{%
  \begingroup
  \renewcommand{\ZSF@figBoxOpts}{}%
  \pgfkeys{/zsf/fig, ZSF@figDefaults, width=0.30, #1}%
  \expandafter\ZSF@figOpen\expandafter{\ZSF@figBoxOpts}{#3}%
    \noindent
    \begin{minipage}[c]{\ZSF@figWidthFrac\linewidth}%
      \centering
      \ZSFimage[width=1, height=\ZSF@figHeightLimit]{#2}%
    \end{minipage}%
    \hfill
    \begin{minipage}[c]{\dimexpr\ZSFfigUsableFraction\linewidth
                              -\ZSF@figWidthFrac\linewidth\relax}%
      \raggedright
      #4%
    \end{minipage}%
    \par
  \end{figbox}% chktex 10
  \endgroup
}
\makeatother

%
%
%
%
%
%
%
%
%
%
%
%
%
\tcbset{split/.style={
  sidebyside,
  sidebyside gap=\ZSFspaceM,
  lower separated=false,
  lefthand width=#1\linewidth,
  before lower={\ZSFBoxSetup},
}}
%
%
%
%
%
%
%
%
%
%
%
%
%
\makeatletter
\newcommand{\ZSF@splitAlign}{top}
\tcbset{
  splitalign/.is choice,
  splitalign/top/.code   ={\renewcommand{\ZSF@splitAlign}{top}},
  splitalign/center/.code={\renewcommand{\ZSF@splitAlign}{center}},
  splitalign/bottom/.code={\renewcommand{\ZSF@splitAlign}{bottom}},
}
\tcbset{ZSF@resolvesplit/.code={\tcbset{sidebyside align=\ZSF@splitAlign}}}
\makeatother
% splitbox ist das rahmenlose, titellose Preset dieses Reglers — der mit
% Abstand häufigste Fall, deshalb ein eigener Name statt vier Optionen.
%
% Sie steht auf derselben Form wie jede andere Box (zsfboxshape) und nimmt mit
% Reglern zurück, was sie rahmenlos macht. Ein 'blankest' wäre kürzer, würde
% aber den Rahmen per Skin entfernen statt per Regler: Ein `frame=soft` vom
% Aufrufer liefe danach ins Leere, weil der Skin gar nichts zeichnet — der
% Regler wäre auf dieser einen Box stillschweigend wirkungslos.
\tcbset{preset/split/.style={
  zsfboxshape,
  bodyparskip=inherit,
  frame=none,
  padx=none, pady=none,
  ZSF@surface=quiet,
  split=0.3,
}}
\makeatletter
\NewDocumentEnvironment{splitbox}{O{}}
  {\renewcommand{\ZSF@imgBudget}{\ZSFfigMaxHeight}%
   \begin{ZSF@box}{preset/split}{}{#1}}{\end{ZSF@box}} % chktex 9
\makeatother

% ------------------------------------------------------------
% Aussagen / Verfahren / Listen
% ------------------------------------------------------------
% Die leise Fassung: dezent, linker Balken im Akzent des Tons. Steht auf derselben
% Form wie die Titelboxen und nimmt genau das zurück, was sie rahmenlos macht
% — Rahmen, Ecken und die obere/untere Polsterung.
%
% Ton, Justierung und Eintritt kommen aus dem Vertrag (zsfboxcontract, via
% zsfboxshape); hier steht nur noch der Unterschied. Die Fläche setzt dieser
% Stil bewusst NICHT — sie kommt vom Flächen-Resolver am Listenende, der Rolle
% und Ton zusammensetzt.
\tcbset{preset/quiet/.style={
  zsfboxshape,
  bodyparskip=inherit,
  frame=none,
  borderline west={\ZSFstatementBarWidth}{0pt}{\ZSFtoneAccent},
  padx=bar, pady=tight,
}}

% ------------------------------------------------------------
% ZSFlist — eine Liste, zwei Nummerierungen
% ------------------------------------------------------------
% Nummerierte Schrittfolge und Faktenliste sind dieselbe Sache: „eine Liste in
% einer optionalen leisen Box". Sie unterscheiden sich in Nummerierung und
% Item-Layout — beides Eigenschaften einer Liste, keine verschiedenen Absichten.
%
%   \begin{ZSFlist}[Titel]            Aufzählung, Marke inline
%   \begin{ZSFlist}[Titel][ordered]   Nummerierung, Marke über dem Text
%   \begin{ZSFlist}                   ohne Titel: ganz ohne Box
%
% Das Item-Layout folgt der Listenart, weil es daran hängt: Ein nummerierter
% Schritt trägt Fliesstext und braucht die Marke als eigene Zeile, ein Fakt
% ist ein Einzeiler. Wer beides entkoppeln muss, bekommt einen Regler — aber
% erst, wenn der Fall wirklich auftritt.
\makeatletter
\newif\ifZSF@listOrdered
\newif\ifZSF@listBoxed
\tcbset{ordered/.is if=ZSF@listOrdered, ordered/.default=true}
\newcommand{\ZSF@listOpen}[1]{%
  \ifZSF@listOrdered
    \begin{enumerate}[
      leftmargin=\ZSFlistLeftMarginOrdered, labelsep=\ZSFlistLabelSep,
      itemsep=\ZSFlistItemSep, topsep=#1, parsep=0pt, partopsep=0pt]%
  \else
    \begin{itemize}[
      leftmargin=\ZSFlistLeftMargin, labelsep=\ZSFlistLabelSep,
      label={\textcolor{\ZSFtoneAccent}{\textbullet}},
      itemsep=\ZSFlistItemSep, topsep=#1, parsep=0pt, partopsep=0pt]%
  \fi
}
\newcommand{\ZSF@listClose}{%
  \ifZSF@listOrdered\end{enumerate}\else\end{itemize}\fi % chktex 9
}
%
%
%
%
%
%
%
%
%
%
\pgfkeys{
  /zsf/listplain/.is family, /zsf/listplain,
  ordered/.code={\tcbset{ordered=#1}},
  ordered/.default=true,
  tone/.code={\tcbset{tone=#1}},
  .unknown/.code={% chktex 26
    \PackageError{ZSF}{%
      Regler '\pgfkeyscurrentname' ist auf einer ZSFlist ohne Titel wirkungslos%
    }{%
      Ohne Titel hat die Liste keine Box, auf die ein Box-Regler wirken kann.
      Entweder einen Titel setzen (dann gelten alle Box-Regler) oder die
      Option weglassen. In dieser Form gelten nur: ordered, tone.%
    }%
  },
}
\NewDocumentEnvironment{ZSFlist}{O{} O{}}
  {\ZSF@listOrderedfalse
   \renewcommand{\ZSF@boxOwnKnobs}{,ordered,}%
   \IfBlankTF{#1}%
     {\ZSF@listBoxedfalse\pgfkeys{/zsf/listplain, #2}\ZSF@listOpen{\ZSFspaceXS}}%
     {\ZSF@listBoxedtrue\begin{defbox}[#1][weight=quiet, #2]\ZSF@listOpen{0pt}}}
  {\ZSF@listClose\ifZSF@listBoxed\end{defbox}\fi} % chktex 9
\makeatother
%
%
%
%
%
%
%
%
%
%
%
%
%
%
\makeatletter
% +m/+g: beide Argumente sind lang. Ein Schritt darf mehrere Absätze tragen
% (Erklärung, Displayformel, Schlusssatz) — xparse-Argumente wären sonst kurz,
% wo das frühere \newcommand lang war.
\NewDocumentCommand{\ZSFItem}{+m +g}{%
  \item
  \IfNoValueTF{#2}%
    {#1}%
    {\IfBlankTF{#1}%
       {#2}%
       {\ifZSF@listOrdered
          {\ZSFfontBlockLabel #1}\\[\ZSFlistMarkGap]#2%
        \else
          {\ZSFfontBlockLabel #1}\,---\,#2%
        \fi}}%
}
\makeatother

%
%
%
%
%
%
%
%
%
%
\makeatletter
\newif\ifZSF@goalStarted
\tcbset{preset/goal/.style={
  zsftitlebox,
  align=center,
  pady=normal,
  frame=strong,
  ZSF@surface=plain,
  /utils/exec={\ZSF@goalStartedfalse},
  after upper={\par},
  atomic,
}}
\makeatother
\NewDocumentEnvironment{goalbox}{O{} O{}}
  {\begin{ZSF@box}{preset/goal}{#1}{#2}}{\end{ZSF@box}} % chktex 9

\newtcbox{\GoalConditionBox}{%
  on line, enhanced,
  boxrule=\ZSFgoalConditionRuleWidth,
  colback=\GoalConditionBackColor,
  colframe=\GoalConditionFrameColor,
  borderline={\ZSFgoalConditionBorderlineWidth}{0pt}{\GoalConditionBorderColor},
  arc=\ZSFgoalConditionInnerArc,
  outer arc=\ZSFboxArc,
  left=\ZSFspaceXS, right=\ZSFspaceXS, top=\ZSFspaceXS, bottom=\ZSFspaceXS,
  nobeforeafter,
  tcbox raise base,
}
\newtcbox{\GoalTargetBox}{%
  on line, enhanced,
  boxrule=\ZSFboxRuleWidth,
  colback=\GoalTargetBackColor,
  colframe=\GoalTargetFrameColor,
  arc=\ZSFboxInnerArc,
  outer arc=\ZSFboxArc,
  left=\ZSFspaceXS, right=\ZSFspaceXS, top=\ZSFspaceXS, bottom=\ZSFspaceXS,
  nobeforeafter,
  tcbox raise base,
}

%
%
%
%
%
%
%
%
%
%
%
%
%
%
\makeatletter
\newcommand{\ZSF@goalLink}[1]{%
  \ifZSF@goalStarted\ZSFInterlude{\ZSFgoalLinkSep}\else\par\ZSF@goalStartedtrue\fi
  #1\par
}
\newcommand{\GoalCondition}[1]{\ZSF@goalLink{\GoalConditionBox{\ZSFfontNote #1}}}
\newcommand{\GoalStep}[1]{\ZSF@goalLink{$\displaystyle #1$}}
\newcommand{\GoalTarget}[1]{\ZSF@goalLink{\GoalTargetBox{$\displaystyle #1$}}}
%
%
%
%
%
%
%
%
%
%
%
\NewDocumentCommand{\ZSFDerivationCase}{s m m m}{%
  \GoalCondition{#2}%
  \IfBooleanTF{#1}%
    {\GoalStep{#3 = \GoalTargetBox{$\displaystyle #4$}}}%
    {\GoalStep{#3}\GoalStep{{}= \GoalTargetBox{$\displaystyle #4$}}}%
}
\makeatother
% Inline-Danger-Tag (rotes Pill-Label für Stolperfallen im Fliesstext)
\newtcbox{\ZSFdanger}{
  colback=DangerTagColor, colframe=DangerTagColor, coltext=\DangerTagTextColor,
  on line, boxsep=\ZSFspaceXS,
  left=\ZSFspaceS, right=\ZSFspaceS, top=\ZSFspaceXS, bottom=\ZSFspaceXS,
  boxrule=0pt, arc=\ZSFboxArc, outer arc=\ZSFboxArc,
  fontupper=\ZSFfontDangerTag
}

% ------------------------------------------------------------
% Valuegrid: kompakte (tabularray-)Wertegrids in einer Panel-Box
% ------------------------------------------------------------
% Interner Panel-Rahmen des Valuegrids — kein öffentlicher Baustein.
% Kapitel greifen ausschliesslich über \begin{valuegrid}{n}[Titel] zu.
\makeatletter
% Die Gitterlinien sind ein Regler, keine Eigenschaft des Bausteins: Ob ein
% Wertegitter seine Zellen einzeln umrandet, hängt am Inhalt und wird pro
% Stelle entschieden (Reglertest, styles/README.md). Ein dichtes Zahlengitter
% will beide Linienrichtungen, eine dreispaltige Zuordnungsliste kommt mit
% waagrechten aus, und wo die Spalten selbst sprechen, stört jedes Lineal.
%
% Das Komma steht IM Wert, nicht in der Spec: `grid=none` liefert sonst einen
% leeren Eintrag in der tblr-Optionsliste.
% Die Linienstaerke ist \ZSFboxRuleWidth und nicht tabularrays Vorgabe: Eine
% Linie ist in diesem System EIN Mass — der Boxrahmen zog 0.5pt, das Gitter
% daneben 0.4pt, und entschieden hatte das niemand.
\newcommand{\ZSF@gridSpec}{{wd=\ZSFboxRuleWidth}}
\newcommand{\ZSF@gridLines}{hlines=\ZSF@gridSpec,vlines=\ZSF@gridSpec,}
\tcbset{
  grid/.is choice,
  grid/both/.code      ={\renewcommand{\ZSF@gridLines}{hlines=\ZSF@gridSpec,vlines=\ZSF@gridSpec,}},
  grid/horizontal/.code={\renewcommand{\ZSF@gridLines}{hlines=\ZSF@gridSpec,}},
  grid/none/.code      ={\renewcommand{\ZSF@gridLines}{}},
  % Dieselben zwei Fragen wie an der ZSFtable, unter denselben Namen und mit
  % derselben Antwort (\ZSF@setRows / \ZSF@setColsep in 20_tables). Sie stehen
  % hier als tcolorbox-Schlüssel, weil das valuegrid eine Box ist und seine
  % Optionsliste die der Box IST — die ZSFtable hat dafür ihre eigene Familie.
  % Zwei Aufrufwege, eine Bedeutung; die Werte kommen aus einer Hand.
  rows/.is choice,
  rows/normal/.code={\ZSF@setRows{normal}},
  rows/roomy/.code ={\ZSF@setRows{roomy}},
  rows/tight/.code ={\ZSF@setRows{tight}},
  colsep/.is choice,
  colsep/normal/.code={\ZSF@setColsep{normal}},
  colsep/tight/.code ={\ZSF@setColsep{tight}},
}
% `rows` und `colsep` gehören in die Vorbelegung: Die Register sind sonst das,
% was die zuletzt gesetzte ZSFtable hinterlassen hat — ausserhalb ihrer Gruppe
% also null, und das Gitter stünde ohne jede Polsterung da.
\tcbset{preset/valuegrid/.style={
  zsftitlebox,
  padx=none, pady=tight,
  ZSF@surface=plain,
  frame=strong,
  grid=both,
  rows=normal, colsep=normal,
  atomic,
}}
% `width` ist Pflicht, nicht Kosmetik: Ohne sie setzt tabularray die Spalten
% auf ihre natuerliche Breite, und das Gitter endet mitten in der Box —
% sichtbar, seit die Optionen ueberhaupt ankommen (siehe unten). Eine Tabelle
% fuellt in diesem System immer die Zeilenbreite (20_tables, Kopf).
%
% Zellpolsterung und Zeilenluft kommen aus DENSELBEN Registern wie bei der
% ZSFtable (\ZSF@setColsep / \ZSF@setRows in 20_tables). Vorher standen hier
% Eigenwerte: `rows=roomy` und `colsep=tight` bedeuteten auf den beiden
% Tabellen-Bausteinen verschiedene Dinge, und \ZSFvaluegridRowSep war ein
% zweites Register fuer dieselbe Frage.
\newcommand{\ZSFvaluegridCommon}{
  width=\linewidth,
  colsep=\ZSF@tblColSep,
  rowsep=\ZSFtableCellPadY,
  \ZSF@gridLines
  rows={valign=m}
}
%
%
%
%
%
%
%
%
%
%
%
%
%
%
%
%
%
%
%
\newcommand{\ZSF@vgOpts}{}
\newcommand{\ZSF@vgOpen}[1]{\begin{tblr}{#1}} % chktex 9
\NewDocumentEnvironment{valuegrid}{O{} O{} m}
  {\renewcommand{\ZSF@boxOwnKnobs}{,grid,rows,colsep,}%
   \renewcommand{\ZSF@imgBudget}{\ZSFimageMaxHeight}%
   \begin{ZSF@box}{preset/valuegrid}{#1}{#2}%
   \edef\ZSF@vgOpts{colspec={X[l] *{#3}{X[c]}},\ZSFvaluegridCommon}%
   \expandafter\ZSF@vgOpen\expandafter{\ZSF@vgOpts}}
  {\end{tblr}\end{ZSF@box}} % chktex 9
\makeatother

%
%
%
%
%
%
%
%
%
%
%
%
%
%
%
%
\newcommand{\ZSFBoxesBreakableOn}{%
  \tcbset{zsftitlebox/.append style={breakable},
          preset/quiet/.append style={breakable}}%
}
\newcommand{\ZSFBoxesBreakableOff}{%
  \tcbset{zsftitlebox/.append style={breakable=false},
          preset/quiet/.append style={breakable=false}}%
}
% Hebt die Platzreserven vor Boxen und Überschriften auf, damit ein
% Anhang dicht packt statt früh umzubrechen.
\newcommand{\ZSFBreakReservesOff}{%
  \renewcommand{\BreakIfNotEnoughSpace}{}%
  \renewcommand{\BreakIfNotEnoughSpaceChapter}{}%
  \renewcommand{\BreakIfNotEnoughSpaceSubsubsection}{}%
}

% % ============================================================
% 65_code_style.tex — OPTIONALES Code-Listing-Modul (CodeExpert)
% ============================================================
% Opt-in: in preamble.tex einkommentieren für code-lastige
% Fächer (Informatik). Das Modul lädt seine optionale Abhängigkeit
% listings samt tcolorbox-Listings-Library selbst.
% Farben (ce_*, code_*) aus 40_colors_structure.tex,
% Längen (\ZSFcode*) aus 30_layout_spacing.tex,
% Box-Grundvertrag (zsftitlebox) und Regler aus 60_boxes.tex.
%
% Stellt bereit:
%   - lstlisting-Style "CodeExpert" (dunkler Code-Block)
%   - codebox[Titel][Optionen]  — Container für Code-Snippets; ohne
%                                 vertikale Polsterung, weil lstlisting
%                                 eigene Skips mitbringt
%   - Regler `part`             — whole (Default) | first | mid | last:
%                                 mehrteilige Code-Gruppen mit offenen
%                                 Kanten, zwischen denen die Spalte
%                                 umbrechen darf
% ============================================================

\RequirePackage{listings}
\tcbuselibrary{listings}

% -------------------------
% CodeExpert Listings Style
% -------------------------
% Was NICHT zum Stil gehört: Satzmechanik, die für jedes Listing gilt.
% Alles, was CodeExpert unten ohnehin setzt, stand hier ein zweites Mal —
% andere Rahmenart, andere Skips — und welcher Wert galt, hing daran, ob der
% Autor den Stil angefordert hatte.
\lstset{
    columns=flexible,
    tabsize=2,
}

\lstdefinestyle{CodeExpert}{
    language=Python,
    basicstyle=\ttfamily\linespread{\ZSFcodeLeadingFactor}\color{code_fg},
    numbers=none,
    aboveskip=\ZSFspaceM,
    belowskip=\ZSFspaceM,
    frame=single,
    rulecolor=\color{code_rule},
    framerule=\ZSFcodeFrameRule,
    framexleftmargin=\ZSFcodeXMargin,
    framexrightmargin=\ZSFcodeXMargin,
    framextopmargin=\ZSFspaceM,
    framexbottommargin=\ZSFspaceM,
    xleftmargin=\ZSFcodeXMargin,
    xrightmargin=\ZSFcodeXMargin,
    backgroundcolor=\color{code_bg},
    keywordstyle=\bfseries\color{ce_pink},
    commentstyle=\color{ce_green},
    stringstyle=\color{ce_yellow},
    identifierstyle=\color{code_fg},
    showstringspaces=false,
    breaklines=true,
    breakatwhitespace=true,
    literate={->}{{$\rightarrow$}}2 {ü}{{\"u}}1 {ä}{{\"a}}1 {ö}{{\"o}}1 {Ü}{{\"U}}1 {Ä}{{\"A}}1 {Ö}{{\"O}}1 {—}{{---}}1 {−}{{-}}1 {„}{{\glqq}}1 {“}{{\grqq}}1 {”}{{''}}1 {’}{{'}}1 {–}{{--}}1 {→}{{$\rightarrow$}}1
}

%
%
%
%
%
%
%
\lstset{style=CodeExpert}

% ------------------------------------------------------------
% codebox — Container für Code-Snippets (optional mit Titel)
%   Identisches Erscheinungsbild wie die Titelboxen (zsftitlebox),
%   aber weisser Hintergrund und zero top/bottom padding, da
%   lstlisting eigene above/belowskips mitbringt.
% ------------------------------------------------------------
\tcbset{preset/code/.style={
  zsftitlebox,
  ZSF@surface=plain,
  pady=none,
}}

%
%
%
%
%
%
%
%
%
%
%
%
%
%
%
%
%
%
%
%
%
%
%
%
%
%
%
%
\tcbset{zsf@codebox@split@base/.style={
  zsfboxcontract,
  unbreakable,
  bodyparskip=inherit,
  ZSF@surface=plain,
  leftrule=\ZSFboxRuleWidth, rightrule=\ZSFboxRuleWidth,
  toprule=0pt, bottomrule=0pt,
  arc=0pt, outer arc=0pt,
  boxsep=0pt,
  padx=normal, pady=none,
  % Der Vorabstand wird über das REGISTER genullt, nicht über `before skip`.
  % Letzteres ist in tcolorbox ein Stil, der `before` schreibt, und überschriebe
  % damit den Eintritts-Hook aus zsfbox (\ZSF@blockBefore): Der Abstand käme an,
  % die Balken-Merker und das \par gingen verloren. Über /utils/exec gilt der
  % Wert nur für diese Box — nachgemessen, der Hook bleibt.
  /utils/exec={\renewcommand{\ZSFBoxBeforeSkip}{0pt}},
  after skip=\ZSFspaceXS,
}}
% Die vier Kantenformen als Auswahl-Schlüssel. 'whole' ist der Default und
% liefert die geschlossene Box; die drei anderen öffnen je eine Kante.
%
% `part` wählt wie `weight` keinen WERT, sondern einen BASISSTIL — und ein
% Basisstil bringt Vorbelegungen mit (`padx=normal`, `pady=none`, Skips).
% Stünde er wie die übrigen Regler in der Aufrufer-Liste, liefen diese
% Vorbelegungen nach den Wahlen des Aufrufers und überschrieben sie:
% `[padx=bar, part=first]` verlor den Innenabstand, `[part=first, padx=bar]`
% behielt ihn — stumm, weil beide Werte für sich gültig sind. Deshalb liest
% die Umgebung die Wahl unten im Vor-Lauf und setzt die fertige Basis VOR die
% Aufrufer-Optionen, genau wie `ZSF@weightBase` in styles/60_boxes.tex.
\tcbset{
  ZSF@partBase/.is choice,
  ZSF@partBase/whole/.style={preset/code},
  ZSF@partBase/first/.style={
    zsf@codebox@split@base,
    toprule=\ZSFboxRuleWidth,
    arc=\ZSFboxInnerArc, outer arc=\ZSFboxArc,
    sharp corners=south,
    % KEIN `before skip` hier. Es ist in tcolorbox ein Stil, der `before`
    % schreibt, und ueberschriebe damit den Eintritts-Hook aus zsfbox
    % (\ZSF@blockBefore) — der Abstand kaeme an, die Balken-Merker gingen
    % verloren. Der Abstand kommt ohnehin von dort, aus demselben Register.
    zsftitlebar,
  },
  ZSF@partBase/mid/.style={zsf@codebox@split@base},
  ZSF@partBase/last/.style={
    zsf@codebox@split@base,
    bottomrule=\ZSFboxRuleWidth,
    arc=\ZSFboxInnerArc, outer arc=\ZSFboxArc,
    sharp corners=north,
    after skip=\ZSFspaceS,
  },
  % Der öffentliche Schlüssel bleibt als No-op bestehen: Er steht weiterhin in
  % der Aufrufer-Liste, wo tcolorbox ihn sonst als unbekannt meldete. Gelesen
  % wurde er zu diesem Zeitpunkt längst.
  part/.is choice,
  part/whole/.code={},
  part/first/.code={},
  part/mid/.code={},
  part/last/.code={},
}
\makeatletter
\newcommand{\ZSF@codeBase}{ZSF@partBase=whole}
\pgfkeys{
  /zsf/codepart/.is family, /zsf/codepart,
  part/.code={\renewcommand{\ZSF@codeBase}{ZSF@partBase=#1}},
  .unknown/.code={},% chktex 26
}
% Öffnet die Box mit EXPANDIERTER Basis — ohne den Umweg landete der Makroname
% selbst in der Optionsliste (dasselbe Muster wie \ZSF@figOpen in 60_boxes).
\newcommand{\ZSF@codeOpen}[3]{\begin{ZSF@box}{#1}{#2}{#3}} % chktex 9
\NewDocumentEnvironment{codebox}{O{} O{}}
  {\renewcommand{\ZSF@codeBase}{ZSF@partBase=whole}%
   \renewcommand{\ZSF@boxOwnKnobs}{,part,}%
   \pgfkeys{/zsf/codepart, #2}%
   \expandafter\ZSF@codeOpen\expandafter{\ZSF@codeBase}{#1}{#2}}
  {\end{ZSF@box}} % chktex 9
\makeatother
    % opt-in: Informatik
\input{styles/66_index}           % opt-in: Stichwortverzeichnis
\ZSFIndexShowPageNumbertrue       % Abschnitt · (S. Seite)
%
%
\ifdefined\ZSFShowcaseMode
  % ============================================================
% 65_code_style.tex — OPTIONALES Code-Listing-Modul (CodeExpert)
% ============================================================
% Opt-in: in preamble.tex einkommentieren für code-lastige
% Fächer (Informatik). Das Modul lädt seine optionale Abhängigkeit
% listings samt tcolorbox-Listings-Library selbst.
% Farben (ce_*, code_*) aus 40_colors_structure.tex,
% Längen (\ZSFcode*) aus 30_layout_spacing.tex,
% Box-Grundvertrag (zsftitlebox) und Regler aus 60_boxes.tex.
%
% Stellt bereit:
%   - lstlisting-Style "CodeExpert" (dunkler Code-Block)
%   - codebox[Titel][Optionen]  — Container für Code-Snippets; ohne
%                                 vertikale Polsterung, weil lstlisting
%                                 eigene Skips mitbringt
%   - Regler `part`             — whole (Default) | first | mid | last:
%                                 mehrteilige Code-Gruppen mit offenen
%                                 Kanten, zwischen denen die Spalte
%                                 umbrechen darf
% ============================================================

\RequirePackage{listings}
\tcbuselibrary{listings}

% -------------------------
% CodeExpert Listings Style
% -------------------------
% Was NICHT zum Stil gehört: Satzmechanik, die für jedes Listing gilt.
% Alles, was CodeExpert unten ohnehin setzt, stand hier ein zweites Mal —
% andere Rahmenart, andere Skips — und welcher Wert galt, hing daran, ob der
% Autor den Stil angefordert hatte.
\lstset{
    columns=flexible,
    tabsize=2,
}

\lstdefinestyle{CodeExpert}{
    language=Python,
    basicstyle=\ttfamily\linespread{\ZSFcodeLeadingFactor}\color{code_fg},
    numbers=none,
    aboveskip=\ZSFspaceM,
    belowskip=\ZSFspaceM,
    frame=single,
    rulecolor=\color{code_rule},
    framerule=\ZSFcodeFrameRule,
    framexleftmargin=\ZSFcodeXMargin,
    framexrightmargin=\ZSFcodeXMargin,
    framextopmargin=\ZSFspaceM,
    framexbottommargin=\ZSFspaceM,
    xleftmargin=\ZSFcodeXMargin,
    xrightmargin=\ZSFcodeXMargin,
    backgroundcolor=\color{code_bg},
    keywordstyle=\bfseries\color{ce_pink},
    commentstyle=\color{ce_green},
    stringstyle=\color{ce_yellow},
    identifierstyle=\color{code_fg},
    showstringspaces=false,
    breaklines=true,
    breakatwhitespace=true,
    literate={->}{{$\rightarrow$}}2 {ü}{{\"u}}1 {ä}{{\"a}}1 {ö}{{\"o}}1 {Ü}{{\"U}}1 {Ä}{{\"A}}1 {Ö}{{\"O}}1 {—}{{---}}1 {−}{{-}}1 {„}{{\glqq}}1 {“}{{\grqq}}1 {”}{{''}}1 {’}{{'}}1 {–}{{--}}1 {→}{{$\rightarrow$}}1
}

%
%
%
%
%
%
%
\lstset{style=CodeExpert}

% ------------------------------------------------------------
% codebox — Container für Code-Snippets (optional mit Titel)
%   Identisches Erscheinungsbild wie die Titelboxen (zsftitlebox),
%   aber weisser Hintergrund und zero top/bottom padding, da
%   lstlisting eigene above/belowskips mitbringt.
% ------------------------------------------------------------
\tcbset{preset/code/.style={
  zsftitlebox,
  ZSF@surface=plain,
  pady=none,
}}

%
%
%
%
%
%
%
%
%
%
%
%
%
%
%
%
%
%
%
%
%
%
%
%
%
%
%
%
\tcbset{zsf@codebox@split@base/.style={
  zsfboxcontract,
  unbreakable,
  bodyparskip=inherit,
  ZSF@surface=plain,
  leftrule=\ZSFboxRuleWidth, rightrule=\ZSFboxRuleWidth,
  toprule=0pt, bottomrule=0pt,
  arc=0pt, outer arc=0pt,
  boxsep=0pt,
  padx=normal, pady=none,
  % Der Vorabstand wird über das REGISTER genullt, nicht über `before skip`.
  % Letzteres ist in tcolorbox ein Stil, der `before` schreibt, und überschriebe
  % damit den Eintritts-Hook aus zsfbox (\ZSF@blockBefore): Der Abstand käme an,
  % die Balken-Merker und das \par gingen verloren. Über /utils/exec gilt der
  % Wert nur für diese Box — nachgemessen, der Hook bleibt.
  /utils/exec={\renewcommand{\ZSFBoxBeforeSkip}{0pt}},
  after skip=\ZSFspaceXS,
}}
% Die vier Kantenformen als Auswahl-Schlüssel. 'whole' ist der Default und
% liefert die geschlossene Box; die drei anderen öffnen je eine Kante.
%
% `part` wählt wie `weight` keinen WERT, sondern einen BASISSTIL — und ein
% Basisstil bringt Vorbelegungen mit (`padx=normal`, `pady=none`, Skips).
% Stünde er wie die übrigen Regler in der Aufrufer-Liste, liefen diese
% Vorbelegungen nach den Wahlen des Aufrufers und überschrieben sie:
% `[padx=bar, part=first]` verlor den Innenabstand, `[part=first, padx=bar]`
% behielt ihn — stumm, weil beide Werte für sich gültig sind. Deshalb liest
% die Umgebung die Wahl unten im Vor-Lauf und setzt die fertige Basis VOR die
% Aufrufer-Optionen, genau wie `ZSF@weightBase` in styles/60_boxes.tex.
\tcbset{
  ZSF@partBase/.is choice,
  ZSF@partBase/whole/.style={preset/code},
  ZSF@partBase/first/.style={
    zsf@codebox@split@base,
    toprule=\ZSFboxRuleWidth,
    arc=\ZSFboxInnerArc, outer arc=\ZSFboxArc,
    sharp corners=south,
    % KEIN `before skip` hier. Es ist in tcolorbox ein Stil, der `before`
    % schreibt, und ueberschriebe damit den Eintritts-Hook aus zsfbox
    % (\ZSF@blockBefore) — der Abstand kaeme an, die Balken-Merker gingen
    % verloren. Der Abstand kommt ohnehin von dort, aus demselben Register.
    zsftitlebar,
  },
  ZSF@partBase/mid/.style={zsf@codebox@split@base},
  ZSF@partBase/last/.style={
    zsf@codebox@split@base,
    bottomrule=\ZSFboxRuleWidth,
    arc=\ZSFboxInnerArc, outer arc=\ZSFboxArc,
    sharp corners=north,
    after skip=\ZSFspaceS,
  },
  % Der öffentliche Schlüssel bleibt als No-op bestehen: Er steht weiterhin in
  % der Aufrufer-Liste, wo tcolorbox ihn sonst als unbekannt meldete. Gelesen
  % wurde er zu diesem Zeitpunkt längst.
  part/.is choice,
  part/whole/.code={},
  part/first/.code={},
  part/mid/.code={},
  part/last/.code={},
}
\makeatletter
\newcommand{\ZSF@codeBase}{ZSF@partBase=whole}
\pgfkeys{
  /zsf/codepart/.is family, /zsf/codepart,
  part/.code={\renewcommand{\ZSF@codeBase}{ZSF@partBase=#1}},
  .unknown/.code={},% chktex 26
}
% Öffnet die Box mit EXPANDIERTER Basis — ohne den Umweg landete der Makroname
% selbst in der Optionsliste (dasselbe Muster wie \ZSF@figOpen in 60_boxes).
\newcommand{\ZSF@codeOpen}[3]{\begin{ZSF@box}{#1}{#2}{#3}} % chktex 9
\NewDocumentEnvironment{codebox}{O{} O{}}
  {\renewcommand{\ZSF@codeBase}{ZSF@partBase=whole}%
   \renewcommand{\ZSF@boxOwnKnobs}{,part,}%
   \pgfkeys{/zsf/codepart, #2}%
   \expandafter\ZSF@codeOpen\expandafter{\ZSF@codeBase}{#1}{#2}}
  {\end{ZSF@box}} % chktex 9
\makeatother

  \input{styles/67_code_comments}
\fi
\input{styles/75_pdf_identity}
\input{styles/70_document_settings}
